\documentclass[10pt,letterpaper]{article}
\usepackage[T1]{fontenc}
\usepackage{lmodern}
\usepackage[margin=0.95in,headheight=15pt]{geometry}
\usepackage{amsmath,amssymb,amsthm,mathtools}
\usepackage{microtype}
\usepackage{booktabs,longtable,array}
\usepackage{enumitem}
\usepackage{xcolor}
\usepackage{fancyhdr}
\usepackage{listings}
\usepackage{xurl}
\usepackage[colorlinks=true,linkcolor=blue!45!black,urlcolor=blue!55!black,
  citecolor=blue!45!black,pdftitle={Formal languages in intuitionistic primitive recursive arithmetic},
  pdfauthor={Research packet prepared for Vladimir Reshetnikov}]{hyperref}
\newcommand{\N}{\mathbb N}
\newcommand{\eps}{\varepsilon}
\newcommand{\ipra}{\ensuremath{\mathsf{iPRA}}}
\newcommand{\PR}{\ensuremath{\mathsf{PR}}}
\newcommand{\Lang}{\mathcal L}
\newtheorem{theorem}{Theorem}[section]
\newtheorem{lemma}[theorem]{Lemma}
\newtheorem{proposition}[theorem]{Proposition}
\newtheorem{corollary}[theorem]{Corollary}
\theoremstyle{definition}
\newtheorem{definition}[theorem]{Definition}
\newtheorem{example}[theorem]{Example}
\theoremstyle{remark}
\newtheorem{remark}[theorem]{Remark}
\setlist{itemsep=2pt,topsep=5pt}
\setlength{\parskip}{3.2pt}
\setlength{\parindent}{0pt}
\setlength{\emergencystretch}{2em}
\pagestyle{fancy}
\fancyhf{}
\fancyhead[L]{\small\textsc{Formal languages in iPRA}}
\fancyhead[R]{\small Constructive proof package}
\fancyfoot[C]{\thepage}
\renewcommand{\headrulewidth}{0.3pt}
\fancypagestyle{plain}[fancy]{}
\lstset{basicstyle=\ttfamily\footnotesize,breaklines=true,
  columns=fullflexible,keepspaces=true,frame=single,
  rulecolor=\color{black!20},showstringspaces=false}

\begin{document}
\begin{titlepage}
\vspace*{0.7in}
{\color{blue!40!black}\Large\textsc{Constructive mathematics / formal languages}}\par
\vspace{0.35in}
{\Huge\bfseries Formal languages in\par
intuitionistic primitive\par
recursive arithmetic}\par
\vspace{0.3in}
{\Large Finite constructions, explicit witnesses,\par
and a Lean formalization blueprint}\par
\vspace{0.45in}
{\large Prepared for Vladimir Reshetnikov}\par
{22 November 2026}\par
\vspace{0.5in}
\begin{minipage}{0.94\textwidth}
\textbf{Abstract.}
We reconstruct a broad part of formal-language theory in first-order
intuitionistic primitive recursive arithmetic. Words, finite-state machines,
grammars, trees, and computations have numerical codes. Positive results
are given by explicit primitive-recursive transformations and bounded
invariants; negative results have concrete separating or diagonal witnesses.
The development includes word algebra, effective Kleene algebra, the
regular-language equivalences and decision procedures, finite monoid
recognition, context-free normalization and parsing, pushdown equivalence,
pumping and nonclosure, noncontracting grammar membership, recognizable
languages and halting obstructions, and exact counting results.
Classical statements that need additional representation or termination
data are reformulated with those data visible. The package contains full
human-readable arguments, an explicit route to object-level Lean
derivations, and independently cross-checked reference algorithms.

\vspace{0.2in}
\textbf{Formal status.}
The Lean material is an uncompiled blueprint. Completed kernel-checked
\texttt{Prf PRAAx} derivations of this report's language theorems are not
claimed. Executable finite checks and document compilation are reported
separately in the verification ledger.
\end{minipage}
\vfill
{\small Repository reference: \texttt{VladimirReshetnikov/ProveIt}\par
Snapshot: \texttt{3eb195cab4cc3d552a9eca78880c6d428f3e44f9}}
\end{titlepage}

\begingroup
\setlength{\parskip}{0pt}
\small
\tableofcontents
\endgroup
\clearpage
\section{The precise constructive setting}
\label{sec:foundations}

\subsection{Scope, representations, and what is being proved}

We work with finite words over a supplied finite alphabet. The alphabet may
vary as numerical input; no assertion here quantifies over a set of arbitrary
sets of words. A language is represented by a finite automaton, a grammar,
a fixed primitive-recursive membership predicate, or a finite-certificate
predicate, as stated in each theorem. Equality of two presentations means
pointwise equivalence on every valid word code. This is a first-order
arithmetical assertion, not equality of objects in a power-set universe.

The target article and its immediately linked mathematical articles supply
the questions, rather than axioms for their answers. The main developments
are the word algebra, regular languages, context-free grammars, bounded
noncontracting derivations, and recognizable languages. Historical and
linguistic claims, classifications of natural languages, and assertions about
arbitrary infinite alphabets are outside the stated finite coding.

\paragraph{Verification status.}
The mathematical proofs below are human-readable derivations with numerical
algorithms, explicit bounds, and quantifier-free induction invariants. The
Lean files supplied with this report are a formalization blueprint; they
have not been compiled in this session. In particular this package does not
contain a completed Lean construction of every object-theory derivation.
Section~\ref{sec:lean} identifies that work precisely. Reference-program
checks are a separate, finite check on the constructions, not a replacement
for a proof or a claim of kernel verification.

\subsection{Language, axioms, and rules of iPRA}

The numerical language has $0$, successor $S$, and, for every \emph{fixed}
primitive-recursive definition, its function symbol. Definitions begin with
zero, successor, and projections, and are closed under composition and
primitive recursion. A recursive definition has equations
\[
 f(0,\vec x)=g(\vec x),\qquad
 f(Sn,\vec x)=h(n,f(n,\vec x),\vec x).
\]
The nonlogical axioms are these defining equations, $Sx\ne0$, successor
injectivity, and induction for quantifier-free formulas, with free numerical
parameters. In sequent form the induction rule is
\[
 \frac{\Gamma\vdash A(0)\quad
       \Gamma,A(n)\vdash A(Sn)}{\Gamma\vdash A(t)},
 \qquad n\text{ fresh for }\Gamma,\quad A\text{ quantifier-free}.
\]
The logical rules are intuitionistic natural deduction for
$\bot,\land,\lor,\to,\forall,\exists$, with equality reflexivity and
substitution. There is no excluded-middle axiom, choice axiom, unrestricted
induction, or rule converting semantic truth into a proof. Negation is
$A\to\bot$. Primitive-recursive predicates abbreviate equations between
their characteristic bits and $1$.

One may add a fixed primitive-recursive definition by a definitional
extension. This does not license a single total primitive-recursive evaluator
for all primitive-recursive programs. Finite automata tables and bounded
machine simulations do have uniform primitive-recursive evaluators; these
are different statements.

\begin{lemma}[Finite arithmetic and coding]\label{base:coding}
Addition, multiplication, truncated subtraction, comparison, exponentiation,
division with a positive divisor, bounded sums and products, and bounded
least search have primitive-recursive definitions whose ordinary defining
properties are provable in \ipra. Finite lists and rectangular tables of
numbers can be encoded with primitive-recursive constructors, selectors,
length tests, and update functions.
\end{lemma}
\begin{proof}
Addition and multiplication use their usual recursions on the second
argument. Predecessor gives truncated subtraction; comparing the two
truncated differences gives equality and order bits. Exponentiation is
repeated multiplication. For division by $d>0$, maintain a quotient--remainder
pair, increasing the remainder unless it reaches $d$, when it is reset to
zero and the quotient increases. Induct on the dividend with invariant
$n=dq+r$ and $r<d$. Bounded sums and products are primitive recursions.
To search below $B$, scan $0,\ldots,B-1$, retaining the first successful
index and otherwise a distinguished failure flag. The invariant that the
retained index is successful and earlier indices fail is a finite Boolean
check of the scanned prefix.

For general list coding use the injective pairing
$\langle x,y\rangle=(x+y)^2+x$ and
$\operatorname{cons}(x,s)=1+\langle x,s\rangle$, with empty list $0$.
The largest $t\le z$ with $t^2\le z$ is bounded-search computable.
On a pairing code, it equals $x+y$, so subtracting its square recovers $x$
and then $y$. Include the range test to reject the gaps in this pairing.
Every proper tail has smaller numerical code. A loop of at most the original
code therefore decodes a list or reports failure. Each operation just
described is a bounded scan or primitive recursion, and each asserted
identity is proved using a quantifier-free invariant of the coded scan.
For entries $<b$ in a table of $m$ cells, the simpler base-$b$ code
$\sum_{i<m}a_i b^i<b^m$ is available when $b\ge2$.
\end{proof}

\begin{lemma}[Boolean compilation and bounded induction]\label{base:boolean}
For every fixed formula $A(\vec x)$ built from primitive-recursive equations
using propositional connectives and bounded numerical quantifiers, there is
a primitive-recursive bit $[A](\vec x)$ such that \ipra{} proves
\[
 A(\vec x)\ \longleftrightarrow\ [A](\vec x)=1,
 \qquad [A](\vec x)=0\ \lor\ [A](\vec x)=1.
\]
Consequently such formulas are decidable and admit induction derived from
quantifier-free induction. A bounded existential has a primitive-recursive
least witness, with a failure value when no witness exists.
\end{lemma}
\begin{proof}
Equality has its comparison bit. On bits, conjunction is product,
disjunction is the sign of the sum, negation is $1\mathbin{\dot-}b$, and
implication is $(1\mathbin{\dot-}b)\lor c$.
For bounded universal quantification use the product of the body bits;
for bounded existential quantification use the sign of their sum.
The finite-scan equations prove that all these expressions are bits.

For clarity, reflection of a bounded quantifier does not use induction on
that quantified formula. The search loop has a bit recording whether a
counterexample has been seen and an index recording its first position.
Its invariant is a conjunction of equations between primitive-recursive
terms, including the product of the tested prefix. If the all-bit is $1$,
instantiate the arbitrary input index below the bound and use the prefix
product to obtain the corresponding body bit. If the all-bit is $0$, the
recorded first counterexample has body bit $0$. Conversely, the hypothesis
that every body holds contradicts this explicit counterexample. The
existential case uses the first successful index in the same way.
Structural recursion on the \emph{fixed formula} now constructs the claimed
intuitionistic derivations. Finally replace $A(n)$ in an induction argument
by the equation $[A](n)=1$, apply quantifier-free induction, and translate
back. Strong bounded induction is obtained by applying this procedure to
the conjunction of all earlier instances.
\end{proof}

\begin{lemma}[Bounded certificate selection]\label{base:selection}
Let $V(x,c)$ be a primitive-recursive bit and $B(x)$ a
primitive-recursive bound. The predicate
$\exists c<B(x)\,V(x,c)=1$ has a primitive-recursive characteristic
function and a primitive-recursive selecting function. Both its success
and failure specifications are provable in \ipra.
\end{lemma}
\begin{proof}
Apply the preceding lemma to the bounded existential and choose the first
successful candidate. The failure value is $B(x)$. Its failure specification
states that each smaller candidate has bit $0$, a bounded formula. No
assertion is made about searching an unbounded set for a witness.
\end{proof}

\begin{lemma}[Finite collection with a supplied bound]\label{base:collection}
For a fixed primitive-recursive $R$ and a supplied numerical bound $B$,
\[
 (\forall i<k\ \exists c<B\ R(i,c))
 \quad\longleftrightarrow\quad
 \exists s<(B+1)^k\ \forall i<k\ R(i,\operatorname{digit}_{B+1}(s,i)),
\]
where the right side also requires every selected digit to be $<B$.
The selecting list is primitive recursive.
\end{lemma}
\begin{proof}
For each $i<k$ select the least witness below $B$ and insert it in the
$i$th base-$(B+1)$ cell. A loop on $i$ builds the table, whose range and
entry conditions are bounded checks. The converse reads the $i$th digit.
The condition remains correct for $B=0$ or $k=0$, using the empty table in
the latter case. In particular this proof does \emph{not} establish
$\forall i<k\exists c\,R(i,c)\Rightarrow\exists s\forall i<k\,R(i,s_i)$
with no common bound. Certificate lists for recognizable-language star
and substitution will be supplied as witnesses in their own right.
\end{proof}

\begin{lemma}[Pigeonhole witnesses and finite graph reachability]
\label{base:graph}
Among $r+1$ numbers below $r$ a repeated pair can be found by bounded
search. In a directed graph on $r$ vertices, every reachable pair has a
path of at most $r-1$ edges (a zero-edge path is allowed). A nonempty
cycle, when present, has a cycle of at most $r$ edges. Reachability,
including shortest-path witnesses, is primitive-recursively decidable.
\end{lemma}
\begin{proof}
For the first claim maintain the $r$ occurrence counts. If every pair of
positions differed, each count would be at most one, whereas the sum is
$r+1$. The two equal positions are then obtained by a scan over fewer than
$(r+1)^2$ pairs. The count equalities and inequalities are bounded-induction
invariants. For a given path with a repeated vertex, delete the segment
between equal occurrences. This preserves its endpoints and each remaining
edge. Repeat for at most its original length; the remaining path has no
repeated vertices. For a nonempty cycle stop upon obtaining a simple cycle
rather than deleting its entire traversal.

Alternatively compute $R_0(i,j)=[i=j]$ and
$R_{t+1}(i,j)=R_t(i,j)\lor\bigvee_{u<r}(R_t(i,u)\land E(u,j))$.
The invariant says that the entry has a path witness of at most $t$ edges
and that every such path makes the entry true. Keep a first predecessor
when an entry becomes true. After $r$ rounds cycle deletion proves
completeness. All loops and all stored path lengths have the stated bounds.
\end{proof}

\begin{lemma}[Finite saturation and bounded structural recursion]
\label{base:saturation}
An inflationary operation on subsets of an $m$-element set stabilizes after
at most $m$ strict changes, provided an unchanged stage is a fixed point.
Finite syntax trees and finite derivation traces admit primitive-recursive
bottom-up and top-down passes with supplied primitive-recursive size bounds.
The usual table invariants are provable in \ipra.
\end{lemma}
\begin{proof}
Represent a subset by $m$ bits. At a changing stage its cardinality increases
by at least one and never exceeds $m$. Induction on the stage, using the
cardinality as a bounded rank, proves stabilization. For a tree, use a
numbering in which children precede their parent, check that ordering, and
scan the node table. At stage $t$, a bounded conjunction states the desired
property of all nodes below $t$. A postorder table can itself be constructed
by bounded scans of a supplied tree code. If outputs grow, use the explicit
majorant from the relevant construction to bound all entries. The claim is
an implementation of these specified passes, not a theorem that every
terminating program has a primitive-recursive running time.
\end{proof}

\begin{proposition}[No universal primitive-recursive evaluator]
\label{base:noeval}
There is no primitive-recursive function $U(e,x)$ whose sections include
every unary primitive-recursive function with their correct values.
\end{proposition}
\begin{proof}
If such a fixed $U$ existed, $d(x)=U(x,x)+1$ would be primitive recursive.
Let $e$ be an index for $d$. Then $d(e)=U(e,e)=d(e)+1$, contradicting
successor arithmetic. The contradiction is constructive. This is an
external impossibility statement about the proposed evaluator; its
diagonal equation is already an \ipra{} argument for each supplied
putative definition and index. It does not prohibit uniform evaluation
of a restricted finite syntax with a primitive-recursive bound, such as
automata or a machine computation supplied with a step counter.
\end{proof}

\subsection{How the later proofs lower to the object theory}

Each finite construction is specified by numerical loops. Each correctness
proof either compares two computed codes, transforms an explicitly supplied
certificate into another one, or checks a finite table whose entries have
known numerical bounds. Induction is on a loop counter, a word length,
a supplied tree size, a supplied trace length, or a decreasing rank with
explicit fuel. Statements about all cells of a current table are Boolean
compiled before induction. An unbounded universal quantifier over input
words is introduced only after proving the result for an arbitrary input.
An existential certificate is eliminated by fixing the supplied code;
induction may then run to its numerical size, but never assumes a bound
on all possible certificates.

For example, equality of two automata is proved by one finite comparison
table and an arbitrary-word run invariant. It is not proved by checking
finitely many sample words and assuming the pattern continues. Likewise,
the grammar parser gets completeness from normalization of an arbitrary
supplied parse tree and an explicit bound on a normalized tree. The absence
of a successful bounded parse then excludes every supplied parse by that
normalization theorem.

\subsection{Repository interface used in the blueprint}

The inspected ProveIt revision is
\texttt{3eb195cab4cc3d552a9eca78880c6d428f3e44f9}.
Its ordinary target is \texttt{LeanProofs.ConciseIPRA.Prf PRAAx}, with
first-order terms \texttt{PTm} and formulas \texttt{Fm PTm}.
The relevant source files were read at this revision; the repository was
not built in this session. The inspected interfaces are
\begin{itemize}
\item \texttt{IPRADerivations.Core.Pairing}: pairing projections and
simultaneous recursion;
\item \texttt{IPRADerivations.Core.SearchParameters}: finite-parameter
bounded least search and success/failure theorems;
\item \texttt{IPRADerivations.Core.Pigeonhole}: explicit repeated-position
witnesses;
\item \texttt{IPRADerivations.Foundations.BoundedFormulas}: fixed-formula
compilation, \texttt{bf\_reflect}, \texttt{bf\_dec}, and
\texttt{bf\_choice};
\item \texttt{IPRADerivations.Foundations.BoundedFixpoint}: a fixed-point
theorem from a bounded ascent invariant;
\item \texttt{IPRABridge.Standard}: transport of closed ordinary iPRA
derivations to the Foundation interface.
\end{itemize}
These are implementation dependencies for future formalization, not a
claim that an imported theorem already proves our language results. The
earlier complexity research packet is useful for compatible bounded-machine
conventions; its semantic Lean results and its proposed object-theory
derivations are kept distinct here as in that packet
\cite{ProveIt,ProveItComplexity}.

\section{Words and effective language algebra}
\label{sec:words}

\subsection{A numerical free monoid}

Fix $q\ge1$ and write $\Sigma_q=\{0,\ldots,q-1\}$. Put $b=q+1$.
Encode a word by
\[
 \operatorname{code}(\eps)=0,\qquad
 \operatorname{code}(wa)=b\operatorname{code}(w)+(a+1).
\]
Thus a valid nonzero code has only digits $1,\ldots,q$ in base $b$;
the most significant digit is its first letter plus one. Dividing repeatedly
by $b$ recovers the word from right to left. The loop needs at most the
original code plus one iterations. Invalid numerical inputs are assigned
a failure flag; every assertion about words includes the validity test.
For $q=0$ take the empty word as the only word and omit letter operations.
This exceptional case satisfies all compatible statements below directly.

\begin{lemma}[Coding, concatenation, and length]\label{word:code}
Word validity, length, letter lookup, concatenation, prefix, suffix,
substring, and reversal are primitive recursive. Coding is injective, and
\[
 |uv|=|u|+|v|,\qquad
 \operatorname{code}(uv)=b^{|v|}\operatorname{code}(u)+
                          \operatorname{code}(v),\qquad
 0\le\operatorname{code}(w)<b^{|w|}.
\]
The last strict bound includes $w=\eps$, since $0<1$.
\end{lemma}
\begin{proof}
Division with remainder uniquely recovers the last digit and shorter prefix.
Induction on the digit count proves injectivity. The concatenation formula
follows by appending one letter at a time; the length equation follows in
the same induction. The bound is preserved by
$bc+(a+1)\le b(b^n-1)+(b-1)=b^{n+1}-1$.
To look up a letter, divide by the appropriate power of $b$ and take the
remainder. Prefixes and suffixes use division and remainder by $b^k$;
substrings are compositions of these operations. Reversal reads the letters
in decreasing position order and appends them to a fresh word. All these
loops are bounded by a supplied length or code, and their digit invariants
are bounded formulas.
\end{proof}

\begin{theorem}[Free-monoid equations and cancellation]\label{word:monoid}
Concatenation is associative, $\eps$ is its two-sided identity, and
\[
 xu=xv\Rightarrow u=v,\qquad ux=vx\Rightarrow u=v.
\]
Moreover $uv=\eps$ exactly when $u=v=\eps$.
\end{theorem}
\begin{proof}
The digit at each position of $(uv)w$ and $u(vw)$ is the same: its position
lies in the $u$, $v$, or $w$ interval in either grouping. Their lengths
agree, so injectivity proves associativity. An empty interval adds no
letters, proving the unit laws. In a cancellation equation, equality of
lengths first gives equal lengths of the uncancelled words; compare the
positions after the common prefix, or before the common suffix. Finally,
the sum of two natural lengths is zero precisely when both are zero.
\end{proof}

\begin{lemma}[Prefix comparison and Levi's lemma]\label{word:levi}
If $uv=xy$, either $|u|\le|x|$ and there is a unique $z$ with
$x=uz$ and $v=zy$, or $|x|\le|u|$ and there is a unique $z$ with
$u=xz$ and $y=zv$. At equality the two descriptions use $z=\eps$.
The applicable factorization is found by a primitive-recursive procedure.
\end{lemma}
\begin{proof}
Compare the two lengths, a decidable numerical comparison. In the first
case the first $|u|$ positions of $x$ equal $u$, since both are prefixes
of the displayed common word. Let $z$ be the remaining suffix of $x$.
Substitute $x=uz$ and cancel $u$. The other case is symmetric.
Uniqueness follows by cancellation; the suffix operation computes $z$.
\end{proof}

\begin{corollary}[Deciding elementary word equations]\label{word:equations}
The equation $ux=v$ has a solution exactly when $u$ is a prefix of $v$;
the solution is then unique and is the suffix after that prefix. The
analogous right-suffix equation is equally decidable. All factorizations
$w=uv$ are exactly the $|w|+1$ cuts at positions $0,\ldots,|w|$.
\end{corollary}
\begin{proof}
Apply the length and position equations. A solution must have length
$|v|-|u|$ and exactly the specified letters. Every cut plainly supplies
a factorization, and its left-factor length uniquely determines it.
\end{proof}

\begin{lemma}[Powers and reversal]\label{word:powers}
The power function $w^0=\eps$, $w^{n+1}=w^nw$ is primitive recursive,
and
\[
 w^{m+n}=w^mw^n,\quad |w^n|=n|w|,\quad
 (uv)^R=v^Ru^R,\quad (w^R)^R=w,\quad (w^n)^R=(w^R)^n.
\]
\end{lemma}
\begin{proof}
The first two identities are induction on the exponent with the monoid
equations. Reversal maps position $i$ to $|w|-1-i$; applying that map
twice is the identity, and its action on the two intervals of $uv$
gives the reversed concatenation. The last identity follows by induction
using this rule. Every induction compares coded words or numerical terms.
\end{proof}

\begin{theorem}[Commuting words are powers of a common word]
\label{word:commute}
For supplied words $u,v$, the equation $uv=vu$ holds exactly when there
are a word $t$ and numbers $r,s$ with $u=t^r$ and $v=t^s$.
When the equation holds, such witnesses are primitive-recursively
computable; one may take $|t|\le |u|+|v|$ and $r,s\le |u|+|v|+1$.
\end{theorem}
\begin{proof}
Common powers commute by the power equations. Conversely, if one word is
empty, use the other as $t$ and exponents $0,1$, taking $t=\eps$ if both
are empty. If both are nonempty and $|u|\le|v|$, Levi's lemma gives
$v=uz$. Cancelling the initial $u$ in $uuz=uzu$ gives $uz=zu$.
The total length of the new pair $(u,z)$ is $|v|$, strictly less than
the previous total. Apply the same procedure to this smaller pair. If
$u=t^r$ and $z=t^s$, then $v=t^{r+s}$, which reconstructs the required
answer. The other length order interchanges the words.

Implement the descent with fuel $|u|+|v|+1$ and store the substitutions
in a finite stack. Every nonterminal stage decreases the positive length
sum. The reverse pass verifies the equations for the stored pair at each
stage. For a nonempty $t$ the exponents are bounded by the corresponding
word lengths; the empty cases have the displayed bounds. Thus the proof
uses bounded descent and a quantifier-free reconstruction invariant, not
an unbounded well-founded recursion.
\end{proof}

\begin{theorem}[Counting and enumerating words]\label{word:count}
There is a primitive-recursive bijection between words of length $n$ over
$\Sigma_q$ and numbers below $q^n$. There are
$\sum_{i=0}^{N}q^i$ words of length at most $N$. If $q\ge1$, the words
$a^n$ for a fixed letter $a$ give arbitrarily many pairwise distinct words.
\end{theorem}
\begin{proof}
For $q\ge2$, read the letters themselves as base-$q$ digits, retaining
the supplied length so leading zeroes are not lost. Successive division
gives the inverse and checks both identities. For $q=1$ there is one word
of every length and one number below $1^n$. Enumerating length blocks
gives the second formula. Distinct exponents give distinct lengths in
the last assertion. This is a numerical enumeration statement; it makes
no cardinality assertion about a power set of languages.
\end{proof}

\begin{theorem}[Extension of a letter map]\label{word:hom}
A finite table $h:\Sigma_q\to\Sigma_r^*$ extends uniquely to a
concatenation-preserving map on words with $h(\eps)=\eps$.
The extension is primitive recursive. With
$M=\max_{a<q}|h(a)|$, it satisfies $|h(w)|\le M|w|$.
If every image is nonempty, then $|w|\le|h(w)|$.
\end{theorem}
\begin{proof}
Read the input from left to right and append the table entry for each
letter. Splitting this scan at a cut proves $h(uv)=h(u)h(v)$.
Any other such map must have the same value after every prefix, by
induction on its length. The length is the sum of the image lengths,
giving both bounds. The same argument proves the universal property
for a represented monoid with primitive-recursive multiplication and
proved identity and associativity laws, replacing concatenation by that
multiplication. This is a schema for each supplied represented monoid,
not a quantifier over arbitrary mathematical structures.
\end{proof}

\subsection{Primitive-recursive languages and their operations}

In this subsection $L(w)$ means $\ell(w)=1$ for a \emph{fixed}
primitive-recursive bit $\ell$, on valid words. Parameters of $\ell$ may
include numerical finite descriptions. $K=L$ means pointwise equivalence.
Put $0=\varnothing$, $1=\{\eps\}$, and $\top=\Sigma_q^*$, with the
context distinguishing the language $1$ from the number $1$.

Concatenation is the bounded predicate
\[
 (KL)(w)\ \longleftrightarrow\
 \exists i\le|w|\;
 K(w[0:i])\land L(w[i:|w|]).
\]
For nonempty $w$, membership in $L^*$ means that some choice of cuts
between its letters makes every nonempty factor belong to $L$;
there are at most $2^{|w|}$ choices. Set $L^*(\eps)=1$.
A mask of internal cuts, or an increasing sequence
$0=i_0<i_1<\cdots<i_k=|w|$, supplies the witness, with
$k\le|w|$. Empty factors can always be discarded.

\begin{theorem}[Decidable closure and witness bounds]\label{alg:prclosure}
Primitive-recursive languages are closed under union, intersection,
complement, difference, concatenation, star, reversal, and inverse word
homomorphism. A nonerasing homomorphic image of a primitive-recursive
language is primitive recursive as well.
\end{theorem}
\begin{proof}
Boolean operations use the membership bits. Concatenation checks the
$n+1$ cuts, and star checks the at most $2^n$ masks just defined, where
$n=|w|$. Reversal tests $L(w^R)$, and inverse image tests $L(h(w))$.
For the last assertion, $h(u)=w$ implies $|u|\le n$, so enumerate all
source words of length at most $n$ and test $L(u)$ and equality of the
image. These enumerations are primitive recursive even if exponentially
large. This argument deliberately excludes arbitrary erasing images:
an empty letter image removes the preimage length bound.
\end{proof}

\begin{theorem}[Boolean and concatenation algebra]\label{alg:semiring}
Union and intersection obey their usual commutative, associative,
idempotent, distributive and absorption laws. For the decidable languages
in this subsection, complement obeys both De Morgan laws and double
complement. With union as addition and concatenation as multiplication,
languages form an idempotent semiring:
\[
 (KL)M=K(LM),\quad 1L=L1=L,\quad 0L=L0=0,
\]
\[
 K(L\cup M)=KL\cup KM,\qquad (K\cup L)M=KM\cup LM.
\]
All operations here are monotone in their non-complemented arguments.
\end{theorem}
\begin{proof}
Boolean identities are checked on the finitely many possible input bits,
using their decidability. For concatenation associativity, a witness on
either side is exactly two nondecreasing cut positions
$0\le i\le j\le|w|$ giving a $K$ prefix, an $L$ middle, and an $M$
suffix. Moving parentheses retains these two cuts. Unit and zero use
the unique empty factor and the impossibility of a factor in $0$.
For distributivity, a chosen factor in a union supplies its disjunct;
conversely retain the same cut and inject that disjunct. Monotonicity
retains the witnesses and uses the corresponding membership implication.
Thus none of the identities needs excluded middle for an unbounded formula.
\end{proof}

\begin{theorem}[Star laws and least closure]\label{alg:star}
For primitive-recursive languages,
\[
 0^*=1^*=1,\quad L^*=1\cup LL^*=1\cup L^*L,\quad
 L\subseteq L^*,\quad L^*L^*=L^*,\quad (L^*)^*=L^*.
\]
The operation is monotone. It is the least represented submonoid containing
$L$: if a primitive-recursive language $X$ contains $\eps$ and $L$ and
is closed under concatenation, then $L^*\subseteq X$.
\end{theorem}
\begin{proof}
Use an empty factor list for $\eps$. A nonempty accepted word has a
first and a last nonempty factor; removing either proves the two unfoldings.
Conversely append the specified factor and remove it if empty. A word of
$L$ uses one factor unless it is empty. Joining two factor lists proves
$L^*L^*\subseteq L^*$, and multiplying by $\eps$ proves the reverse.
For $(L^*)^*$, select the inner cut mask for each outer block and flatten
the lists. Every inner block has length at most the original word length,
so the selections have a common bound $2^{|w|}$ and bounded collection
is available. Monotonicity keeps the same mask. For least closure,
multiply the factors in sequence; after $i$ factors the resulting prefix
lies in $X$. This induction uses the bit for $X$ and a supplied finite
factor table. The assertions for $0$ and $1$ follow by discarding empty
factors.
\end{proof}

\begin{proposition}[Positive closure and star identities]\label{alg:identities}
Write $L^+=LL^*$. Then
\[
 L^+=L^*L,\qquad L^*=1\cup L^+,\qquad
 (KL)^*K=K(LK)^*,\qquad
 (K\cup L)^*=K^*(LK^*)^*.
\]
Also $(1\cup L)^*=L^*$ and $L\subseteq K\Rightarrow L^*\subseteq K^*$.
\end{proposition}
\begin{proof}
For $L^+$ keep a list with at least one factor, allowing that factor to be
empty when $\eps\in L$. Separating the first or last factor proves the
first two equations, including that empty-word case. A word in
$(KL)^*K$ is a chain $k_0l_1k_1\cdots l_mk_m$ with each $k_i\in K$
and $l_i\in L$; moving the parentheses proves the sliding identity.
Empty factors cause no change and may be retained in this witness chain.
For the last displayed identity, take a factorization with each factor
tagged $K$ or $L$. Tags can be selected by their membership bits. Combine
the initial run of $K$ factors, and then every $L$ factor with the following
run of $K$ factors. This gives $K^*(LK^*)^*$. Expanding the grouped lists
gives the converse. All selections and regroupings are scans of bounded
factor tables. Adding or removing empty factors proves $(1\cup L)^*=L^*$;
the final monotonicity statement was also proved above.
\end{proof}

\begin{theorem}[Arden's lemma with an effective uniqueness proof]
\label{alg:arden}
For primitive-recursive languages $A,B$, $A^*B$ is a solution of
$X=AX\cup B$ and is contained in every primitive-recursive solution.
If $\eps\notin A$, it is the unique primitive-recursive solution.
\end{theorem}
\begin{proof}
The unfolding law gives $A^*B=B\cup A(A^*B)$. In any solution $X$,
$B\subseteq X$ and $AX\subseteq X$. Induction on the length of a
supplied list of $A$ factors proves $A^*B\subseteq X$.

Assume $\eps\notin A$ and fix $w\in X$. Test $B(w)$. If it holds,
stop. If it fails, the equation supplies $w=uv$ with $u\in A$ and
$v\in X$; choose a witnessing cut by bounded search. Since $|u|\ge1$,
the residual word is strictly shorter. Repeat. After at most $|w|$
positive cuts the residual is empty. At that stage the $AX$ alternative
is impossible, so the residual is in $B$. The stored prefixes are a
factorization in $A^*B$. The invariant records the original word as the
concatenation of the stored $A$ factors and the current residual in $X$;
it is a bounded check. This proves the reverse inclusion without an
unbounded unfolding argument. If $\eps\in A$, uniqueness need not hold:
for $A=1,B=0$, every $X$ solves the equation.
\end{proof}

\begin{theorem}[Reversal of language operations]\label{alg:reverse}
Reversal is an involution on represented languages and satisfies
\[
 (K\cup L)^R=K^R\cup L^R,\quad
 (K\cap L)^R=K^R\cap L^R,\quad
 (\overline L)^R=\overline{L^R},\quad
 (KL)^R=L^RK^R,\quad (L^*)^R=(L^R)^*.
\]
\end{theorem}
\begin{proof}
Membership in $L^R$ is tested at $w^R$ and reversal is an involution,
proving the Boolean equations. Reverse the two factors of a concatenation
and exchange their order. For star reverse every factor and the order of
the factor list; this is a bounded transformation and its own inverse.
\end{proof}

\begin{theorem}[Homomorphic image and inverse-image laws]
\label{alg:homlaws}
Inverse word homomorphism preserves union, intersection, complement, and
star inclusion $h^{-1}(L)^*\subseteq h^{-1}(L^*)$. For any finite letter
map, the existentially represented image satisfies
\[
 h(K\cup L)=h(K)\cup h(L),\quad h(KL)=h(K)h(L),\quad
 h(L^*)=h(L)^*.
\]
In the last equation an image-star witness includes a finite list of source
words. This is the witness-list presentation, with no unbounded collection
step from separate existential preimage assertions. If $h$ is nonerasing,
all these image languages have the bounded,
primitive-recursive readings of Theorem~\ref{alg:prclosure}.
\end{theorem}
\begin{proof}
The inverse-image equations simply substitute $h(w)$ into the membership
formula. If a list factors $w$ into inverse images of $L$, applying $h$
to its factors gives a factorization of $h(w)$ in $L^*$, allowing empty
images. A witness for an image is a supplied source word. For union keep
that word and its disjunct. For concatenation split the source word in the
forward direction; in the reverse direction concatenate the two supplied
source words and use $h(uv)=h(u)h(v)$.
For star apply $h$ to the source factor list, or concatenate the supplied
list of source words in the other direction. Empty image factors may be
deleted on the output side; the source list remains part of the certificate.
No unbounded collection principle is used to create such a list from
pointwise existential statements.
\end{proof}

\begin{proposition}[The source article's arithmetic syntax is regular]
\label{alg:syntax}
Let $D$ be a decimal digit, $D_+$ a nonzero digit, and put
\[
 N=0\ \cup\ D_+D^*,\qquad E=N(\mathtt{+}N)^*,\qquad
 T=E\ \cup\ E\mathtt{=}E.
\]
Then $T$ is exactly the language of numerals, sums, and one equality
described in the source article's example. In particular it has a
primitive-recursive membership test and is regular. It includes
\texttt{23+4=555} regardless of the arithmetic truth of that equation.
\end{proposition}
\begin{proof}
The numeral expression permits exactly $0$ or a nonzero first digit
followed by any digits. Splitting at plus signs gives a unique list of
nonempty numeral fields. Splitting at the optional unique equality sign
gives one or two such sums. These are exactly the permitted recursive
syntax rules: no equality may occur inside either side of the unique
equality, and each plus sign separates valid numeral or sum substrings.
Conversely, repeated splitting of any permitted sum yields those numeral
fields. This splitting decreases the substring length and has at most
the original length many steps. The expression is built from letters by
union, concatenation, and star, hence is a regular expression; the
regular-expression conversion in Section~\ref{sec:regular} supplies a
finite automaton. Evaluating decimal numerals or addition was never part
of this membership condition.
\end{proof}

\section{Regular languages: effective equivalences, closure, and decisions}
\label{sec:regular}

All alphabets in this section are finite, with letters numbered below a supplied
integer $k\geq 1$.  A state set is $[n]=\{0,\ldots,n-1\}$, with $n\geq1$.
A deterministic automaton is a finite transition table
$\delta:[n]\times[k]\to[n]$, an initial state $q_0<n$, and an accepting
bit-vector $F:[n]\to\{0,1\}$.  A nondeterministic automaton replaces
$\delta$ by a Boolean relation $E(q,a,r)$ and permits a bit-vector $I$ of
initial states.  An $\eps$-automaton also has a Boolean relation $E_\eps(q,r)$.
All these data are number codes.  Well-formedness is a bounded test; malformed
codes can be sent to a fixed rejecting automaton.  Consequently every
construction below denotes a total primitive recursive function on codes.

A path is a finite list of states and edge labels.  Its validity, endpoints,
and the word obtained by deleting $\eps$ labels are primitive recursive.
The accepted language is written $L(M)$.  This notation abbreviates an
arithmetic predicate, not an additional sort of arbitrary sets.  In
particular, the equality $L(M)=L(N)$ abbreviates equality of their accepting
bits for every word.  All numerical bounds below are deliberately elementary;
optimal running times are not needed for their interpretation in \ipra.
The corresponding standard Lean interfaces are documented in
\cite{MathlibDFA,MathlibNFA,MathlibEpsilon}; our number-coded implementations
must additionally establish the primitive-recursive and object-theoretic
bridges described below.

\subsection{Finite runs and elimination of nondeterminism}

\begin{lemma}[Run composition]\label{reg:run}
There is a primitive recursive function $\delta^*(q,w)$ satisfying
\[
 \delta^*(q,\eps)=q,\qquad
 \delta^*(q,wa)=\delta(\delta^*(q,w),a),\qquad
 \delta^*(q,uv)=\delta^*(\delta^*(q,u),v).
\]
Thus membership in $L(M)$ is primitive recursive, uniformly in a finite
well-formed DFA code $M$.
\end{lemma}
\begin{proof}
Store the current state and scan the coded input from left to right, taking
one table lookup per letter.  After $i$ letters the stored state is the state
at position $i$ of the unique valid run.  This is an equality between
primitive recursive terms.  Induction on the number of scanned letters
proves the invariant.  Starting the second scan at the state reached by the
first proves the third equation by the same quantifier-free induction.
Acceptance is the final lookup $F(\delta^*(q_0,w))$.
\end{proof}

\paragraph{Finite-graph facts used below.}
The bounded reachability result~\ref{base:graph} applies to every graph
whose vertices are the numbered states of an automaton.  It supplies a
primitive recursive path-shortening operation, a reachability test, and
shortest path witnesses of length below the number of states.  Edge
labels can be retained during each deletion, so a shortened automaton
path gives an explicitly coded shorter word.  The deletion preserves
endpoints, not generally the word itself; where the word must remain
unchanged, we shorten only blocks of $\eps$ edges.

\begin{proposition}[$\eps$-closure has a finite bound]\label{reg:epsilon}
For an $n$-state $\eps$-automaton, the $\eps$-closure $C(S)$ of a state
bit-vector $S$ is primitive recursive.  It is computed by $n$ iterations of
\[
 S\longmapsto S\cup\{r:\exists q<n\;(q\in S\wedge E_\eps(q,r))\}.
\]
It satisfies $S\subseteq C(S)$, $C(C(S))=C(S)$, and preserves finite unions.
If a word of length $\ell$ is accepted, it has an accepting path with at most
$\ell+(\ell+1)(n-1)$ edges.
\end{proposition}
\begin{proof}
After $j$ iterations the bit at $r$ is set exactly when there is an
$\eps$-path of length at most $j$ from $S$ to $r$.  The two directions are
proved by appending an edge or deleting the last edge of a nonempty path.
By the bounded reachability result~\ref{base:graph}, $n-1$ edges already suffice; using $n$ iterations
avoids an exceptional subtraction.  Closure and idempotence follow by
concatenating and shortening paths.  A path from a union starts in one of
its finitely many constituent sets, which proves the union equation.
An accepting path for a length-$\ell$ word has exactly $\ell$ letter edges
and $\ell+1$ intervening $\eps$ blocks.  Shorten each block separately to
length at most $n-1$.  The word and the state at every letter edge are
unchanged, giving the stated bound.  This also replaces the apparently
unbounded existential path definition of acceptance by bounded search.
\end{proof}

\begin{theorem}[Subset construction]\label{reg:subset}
Every $n$-state $\eps$-NFA has an equivalent DFA with $2^n$ states.  The
construction and its correctness proof are primitive recursive and
formalizable in \ipra.  Ordinary NFAs and DFAs consequently define the same
languages as $\eps$-NFAs.
\end{theorem}
\begin{proof}
A DFA state is an $n$-bit subset.  Its initial subset is $C(I)$ and its
letter transition is
\[
 T(S,a)=C\bigl(\{r:\exists q<n\;(q\in S\wedge E(q,a,r))\}\bigr).
\]
A subset accepts when it meets $F$.  For every prefix $u$ of the input,
maintain the invariant that the current bit-vector consists exactly of the
endpoints of paths labelled $u$, allowing a final $\eps$ block.  For the
empty prefix this is the definition of $C(I)$.  For $ua$, split an accepting
path immediately before its last letter edge; conversely concatenate a path
for $u$, that edge, and a final $\eps$ block.  The bounded witnesses from
Proposition~\ref{reg:epsilon} make the invariant a bounded Boolean
assertion for the particular input.  Quantifier-free induction proves it.
A DFA is an NFA with singleton transitions and initial set, and an NFA is
an $\eps$-NFA with no $\eps$ edges.  These embeddings preserve paths.
\end{proof}

\begin{proposition}[Completion and finite languages]\label{reg:finite}
A partial $n$-state DFA can be completed with at most $n+1$ states.  Every
explicit finite list of words $w_1,\ldots,w_t$ defines a regular language,
with a complete DFA using at most $2+\sum_i|w_i|$ states.
\end{proposition}
\begin{proof}
For completion, add a nonaccepting sink, send every missing transition to
it, and loop there on every letter.  A formerly defined run is unchanged;
a run that encounters a missing transition can never accept.
For the second assertion, make states for all distinct prefixes of listed
words, including $\eps$, and add a sink.  From prefix $u$ on letter $a$, go
to $ua$ if it is a listed prefix, otherwise to the sink.  Mark exactly the
listed complete words as accepting.  Induction along a word says that its
run is either at that prefix or at the sink.  There are at most
$1+\sum_i|w_i|$ prefixes; comparing the finitely many codes removes
duplicates constructively.  The empty list and lists containing $\eps$
are handled by the same construction.
\end{proof}

\subsection{Closure constructions}

\begin{theorem}[Boolean closure]\label{reg:boolean}
Regular languages are closed under union, intersection, difference,
symmetric difference, and complement.  Binary operations on $n$- and
$m$-state complete DFAs use $nm$ states; complement uses $n$ states.
\end{theorem}
\begin{proof}
The product automaton has initial state $(q_0,r_0)$ and transition
$(q,r)\mapsto(\delta_M(q,a),\delta_N(r,a))$.  Its run on $w$ is the pair
of the two runs by Lemma~\ref{reg:run}.  For union, intersection,
difference, or symmetric difference, apply the corresponding Boolean
operation to their accepting bits.  Complement replaces $F(q)$ by
$1-F(q)$ and requires no other change.  The complement statement uses
only decidability of a finite accepting bit, not excluded middle for an
arbitrary arithmetic formula.
\end{proof}

\begin{theorem}[Concatenation and star]\label{reg:concatstar}
If $M,N$ have $n,m$ states, there are $\eps$-NFAs for
$L(M)\cup L(N)$, $L(M)L(N)$, and $L(M)^*$ with respectively
$n+m$, $n+m$, and $n+1$ states.  Thus all three languages are regular.
\end{theorem}
\begin{proof}
First rename states so the two copies are disjoint.  For union take the
union of initial sets and the union of final sets, with no edges between
the copies.  Every path stays in its initial copy.
For concatenation use the first initial set, the second final set, and add
$\eps$ edges from every final state of the first copy to every initial
state of the second.  Every accepting path crosses once from the first
copy to the second.  Splitting there gives a factorization $uv$ with
$u\in L(M)$ and $v\in L(N)$; concatenating two witnessing paths gives the
converse.  Empty factors cause no difficulty.
For star add a new state $z$ that is both the sole initial state and the
sole final state.  Add $\eps$ edges from $z$ to every old initial state
and from every old final state to $z$.  A nonempty accepting path splits
at its visits to $z$ into finitely many old accepting paths.  The empty
path represents the zero-factor product.  Conversely concatenate any
finite list of old accepting paths with the connecting $\eps$ edges.
All splits and concatenations act on supplied finite certificates.
Theorem~\ref{reg:subset} then supplies complete DFAs.
\end{proof}

\begin{theorem}[Reversal]\label{reg:reverse}
Regular languages are closed under word reversal.  An $n$-state NFA for
$L$ yields an $n$-state NFA for $L^R$.
\end{theorem}
\begin{proof}
Reverse every edge, interchange the initial and final bit-vectors, and
leave each edge label unchanged.  Reversing a path reverses its sequence
of nonempty labels, and doing it twice returns the original path.
This is a primitive recursive bijection of accepting path certificates.
\end{proof}

\begin{theorem}[Homomorphisms and inverse homomorphisms]\label{reg:morphism}
Let $h$ be specified by a finite table of output words $h(a)$.  Regular
languages are closed under $h$ and $h^{-1}$, including erasing
homomorphisms.  A direct-image construction from an $n$-state DFA uses at
most $n(1+\sum_a|h(a)|)$ $\eps$-NFA states; an inverse-image construction
uses the original number of DFA states.
\end{theorem}
\begin{proof}
For direct image replace every edge labelled $a$ by a private path labelled
$h(a)$.  If $h(a)=\eps$, use an $\eps$ edge; if its length is $d>0$, use
$d$ edges and $d-1$ fresh intermediate states.  Intermediate states belong
to a single replaced edge and have no additional incident edges.  Every
accepting path therefore splits, at the original states, into completed
replacements of original edges.  Its label is $h(u)$ for an accepted
original word $u$.  Conversely replace the edges of any accepting path.
The exact additional state count is
$n\sum_a\max(|h(a)|-1,0)$, bounded as stated.
For inverse image replace the transition on input letter $a$ by
$q\mapsto\delta^*(q,h(a))$.  Run composition gives
$\delta_{\mathrm{new}}^*(q,w)=\delta^*(q,h(w))$, by induction on the
number of letters of $w$.  Initial and final states are unchanged.
\end{proof}

\begin{theorem}[Regular substitution]\label{reg:substitution}
Suppose each source letter $a$ is assigned a regular language $K_a$ by a
finite automaton.  Replacing each letter of words in a regular language
$L$ by an independently chosen word of $K_a$ produces a regular language.
If the source DFA has $n$ states and the automata for the $K_a$ have
$n_a$ states, $n+n\sum_a n_a$ states suffice for an $\eps$-NFA.
\end{theorem}
\begin{proof}
For each source transition $q\mathrel{\smash{\xrightarrow{a}}}r$ take a
private copy of the automaton for $K_a$.  Remove the source edge, put
$\eps$ edges from $q$ to the copy's initial states, and from its accepting
states to $r$.  Keep the source initial and final states.  An accepting
path splits at successive source states into accepting paths through the
private copies, giving a source word and one replacement word per letter.
Conversely, given these finite witnesses, concatenate the corresponding
paths.  A copy can accept $\eps$, but a completed traversal still records
one source letter; no minimum length assumption is used.  There is one
copy per source state and source letter, proving the bound.
\end{proof}

\begin{theorem}[Quotients by a regular language]\label{reg:quotient}
For regular $L,R$, the languages
\[
 L/R=\{u:\exists v\;(v\in R\wedge uv\in L)\},\qquad
 R\backslash L=\{v:\exists u\;(u\in R\wedge uv\in L)\}
\]
are regular.  Given DFAs of $n,m$ states, the existential witnesses needed
to compute the new state sets may always be chosen with length below $nm$.
\end{theorem}
\begin{proof}
For $L/R$, retain the transition table of $L$ and mark $q$ accepting
exactly when a final pair is reachable from $(q,r_0)$ in the product of
$L$ and $R$.  Such a path has length below $nm$ by
the bounded reachability result~\ref{base:graph}.  Therefore its existence is a bounded test, and the
new accepting condition is exactly the displayed definition.
For $R\backslash L$, mark $q$ initial in an NFA exactly when
$(q,r)$ is reachable from $(q_0,r_0)$ for some accepting state $r$ of $R$.
Keep the old transitions and final states of $L$.  The same bound provides
a word $u$ reaching each marked $q$.  A run starting at $q$ on $v$ is
accepting exactly when $uv\in L$.  The two path concatenations prove both
inclusions.  The assertion is intentionally restricted to supplied
regular descriptions of $R$: an arbitrary predicate for $R$ does not
supply a decision procedure for these new finite state sets.
\end{proof}

\begin{corollary}[Prefixes, suffixes, and factors]\label{reg:factors}
The prefixes, suffixes, and contiguous factors of a regular language are
regular.  Their witnessing extension words can be chosen with length
below $n$ on each side, for an $n$-state DFA.
\end{corollary}
\begin{proof}
For prefixes, mark as accepting all states from which a final state is
reachable.  For suffixes, use all states reachable from the original
initial state as NFA initial states.  For factors, make both changes.
A path witnessing an extension can be shortened without changing its
endpoints, so each extension has length below $n$.  Concatenation of the
left extension, the proposed factor, and the right extension proves the
three characterizations.
\end{proof}

\subsection{Regular expressions and regular grammars}

A regular expression is a finite syntax tree generated by
$0,1,a,+,\cdot,{}^*$.  Its denotation interprets these symbols as the empty
language, $\{\eps\}$, a singleton letter, union, concatenation, and star.
A membership certificate records choices at union nodes, word splits at
product nodes, and a finite factor list at star nodes.  Empty factors can
be deleted at a star without changing its product.  For a nonempty word
of length $\ell$, at most $\ell$ remaining factors are needed.  For the
empty word, the zero-factor certificate always suffices.

\begin{proposition}[Bounded regular-expression membership]\label{reg:regexdec}
Membership in a regular expression is primitive recursive uniformly in
the expression and word codes.  It has an algorithm using only a finite
table indexed by expression nodes and input intervals.
\end{proposition}
\begin{proof}
For a fixed word $w$ of length $\ell$, store a Boolean value for every
subexpression $P$ and interval $[i,j)$, $0\leq i\leq j\leq\ell$.
The values for $0,1,a$ are immediate.  Union takes the disjunction.
Product searches the $j-i+1$ splits $i\leq t\leq j$.  For $P^*$ first
set every empty interval true; a nonempty interval is true exactly when
some strictly increasing sequence
$i=t_0<t_1<\cdots<t_r=j$ has every $P$ interval true.  There are at most
$2^{j-i-1}$ such cut sets.  Removing empty factors proves that this
bounded condition is equivalent to star membership.  Evaluate nodes in
postorder, so their children have already been evaluated.  The table has
at most $s(\ell+1)^2$ entries when there are $s$ syntax nodes.  Certificate
extraction stores the first successful split or cut set.  Conversely,
each stored choice assembles a valid certificate.  This proves soundness
and completeness by a bounded conjunction over the table, not by an
induction whose hypothesis quantifies over arbitrary languages or words.
\end{proof}

\begin{theorem}[Regular expressions to automata]\label{reg:thompson}
An expression with $s$ syntax nodes has an equivalent $\eps$-NFA with at
most $2s$ states.  The conversion is primitive recursive.
\end{theorem}
\begin{proof}
Give every subexpression a fresh entry and exit.  Its entry has no incoming
edges inside that subautomaton, and its exit has no outgoing edges there.
For $0$ put no edge between them, for $1$ put an $\eps$ edge, and for a
letter put that labelled edge.  For $P+Q$, connect the new entry to the
two child entries, and the two child exits to the new exit, using $\eps$
edges.  For $PQ$, connect the new entry to $P$'s entry, $P$'s exit to
$Q$'s entry, and $Q$'s exit to the new exit.  For $P^*$ connect the new
entry to the new exit and to $P$'s entry, and $P$'s exit to $P$'s entry
and the new exit.  All these connecting edges are $\eps$ edges.  These
instructions allocate precisely two states per
syntax node.
For union, an entry-to-exit path chooses one child.  For product, it
visits the two child subautomata in order.  For star, split it at returns
to the child entry after the child exit: this gives finitely many child
accepting paths, or the direct zero-factor path.  The converse constructs
paths from membership certificates with the same connectors.  The
entry/exit invariant prevents any other way of moving between children.
For formalization, fix the input word and prove these statements for all
subexpressions and input intervals in the finite table of
Proposition~\ref{reg:regexdec}; shorten $\eps$ blocks when necessary.
Thus the structural proof has an explicitly bounded quantifier-free
invariant.
\end{proof}

\begin{theorem}[Automata to regular expressions]\label{reg:eliminate}
Every finite automaton has an equivalent regular expression, obtained by
a primitive recursive construction.  Together with
Theorem~\ref{reg:thompson}, this proves Kleene's equivalence in \ipra.
\end{theorem}
\begin{proof}
It suffices to treat an $n$-state automaton whose edges carry letters or
$\eps$.  For $i,j<n$ let $R^{(0)}_{ij}$ be the sum of all labels of direct
edges from $i$ to $j$, adding $1$ when $i=j$ and writing $0$ for an empty
sum.  For $0\leq h<n$ define
\[
 R^{(h+1)}_{ij}
 =R^{(h)}_{ij}
  +R^{(h)}_{ih}\bigl(R^{(h)}_{hh}\bigr)^*R^{(h)}_{hj}.
\]
The invariant says that $R^{(h)}_{ij}$ describes exactly the labels of
paths from $i$ to $j$ all of whose internal vertices are below $h$.
At $h=0$ such a path is a single edge or the empty path.  For the induction
step, a path with no internal occurrence of $h$ uses the first summand.
Otherwise split at its first and last internal occurrences of $h$ and at
all the occurrences in between.  The first and last pieces have no
internal $h$, and the middle pieces are $h$-to-$h$ paths with no internal
$h$.  Their labels belong to the three factors of the second summand.
Conversely concatenate these paths; every new internal vertex is either
an old one below $h$ or $h$ itself.  This argument permits $i=h$, $j=h$,
and empty pieces.
At $h=n$ every internal vertex is allowed.  Sum $R^{(n)}_{ij}$ over
initial $i$ and final $j$.  Path splitting and concatenation prove equality
with the accepted language.
The recursion has exactly $n$ rounds and $n^2$ expressions per round.
If the largest base expression has size $b$, the recurrence for sizes is
$s_{h+1}\leq4s_h+4$, hence $s_h\leq4^h(b+2)$.  The final finite sum has
size at most $2n^2(s_n+1)$.  These coarse bounds show directly that even
an implementation that copies syntax trees is primitive recursive.
For each fixed input word the invariant can be expressed over its
finitely many intervals and endpoint pairs; a syntactic path certificate
is shortened block by block, so no appeal to an unbounded search is hidden
in the proof.
\end{proof}

\begin{theorem}[Right-linear grammars]\label{reg:grammar}
Finite right-linear grammars and finite automata define the same languages.
This remains true when productions have forms $A\to uB$ and $A\to u$,
where $u$ is any finite terminal word, including $\eps$.
\end{theorem}
\begin{proof}
First take productions $A\to aB$, $A\to B$, $A\to a$, and
$A\to\eps$.  Make one state for each nonterminal and one fresh accepting
state $f$.  Use the start nonterminal as initial state.  The four
production types give respectively an $a$ edge to $B$, an $\eps$ edge
to $B$, an $a$ edge to $f$, and an $\eps$ edge to $f$.  State $f$ has
no outgoing edges.  Before the last production, a derivable sentential
form is a terminal prefix followed by exactly one nonterminal.  Each
production extends the prefix and updates that state, while the final
production removes it.  This invariant converts a finite derivation
into an accepting path and conversely.
For a general terminal word $u$, replace the production by a private
chain spelling $u$; use an $\eps$ edge if $u$ is empty.  The same invariant
holds at the chain boundaries.
Conversely an NFA gives productions $A_q\to aA_r$ for every letter edge,
$A_q\to A_r$ for every $\eps$ edge, and $A_q\to\eps$ for every final
state.  A fresh start nonterminal has unit productions to initial states.
Again derivations and paths convert by scanning their finite steps.
All added chains and productions have explicit lengths supplied by the
input grammar or automaton, so the conversions are primitive recursive.
\end{proof}

\paragraph{Application of Arden's equation.}
The general PR-language result~\ref{alg:arden} applies in particular to
regular languages.  For $\eps\notin A$, the language $A^*B$ is the unique
PR solution of $X=AX\cup B$.  The closure constructions above additionally
produce its finite automaton from those for $A,B$.  This is useful when
solving finite systems of state equations; the explicit elimination
recurrence in Theorem~\ref{reg:eliminate} does not need to assume that
any particular coefficient is $\eps$-free.

\subsection{Pumping and explicit separation examples}

\begin{theorem}[Regular pumping with computed cuts]\label{reg:pumping}
Let $M$ be a DFA with $n$ states accepting $w$ with $|w|\geq n$.  Primitive
recursive functions compute $x,y,z$ such that
\[
 w=xyz,\qquad |xy|\leq n,\qquad 1\leq|y|,\qquad
 xy^tz\in L(M)\quad\hbox{for every }t\in\N.
\]
\end{theorem}
\begin{proof}
Inspect the first $n+1$ states of the run.  Finite pigeonhole supplies
$i<j\leq n$ with equal states; a bounded search selects the first such
pair.  Put $x=w[0:i]$, $y=w[i:j]$, $z=w[j:|w|]$.  The run on $y$ loops
at the state reached by $x$.  Quantifier-free induction on $t$ proves
that $y^t$ also loops there.  Appending the fixed suffix $z$ reaches the
same accepting state as the original run.  The cut and length assertions
follow from $i<j\leq n$.
\end{proof}

\begin{corollary}[Four explicit nonregular languages]\label{reg:nonregular}
Over letters $a\ne b$, the languages
\[
 \{a^rb^r:r\in\N\},\qquad
 \{w:w=w^R\},\qquad
 \{ww:w\in\{a,b\}^*\}
\]
are not regular.  The unary language $\{a^{r^2}:r\in\N\}$ is also not
regular.  Each assertion is proved in \ipra\ as the refutation of every
purported finite DFA code recognizing that particular PR predicate.
\end{corollary}
\begin{proof}
Suppose an $n$-state DFA recognizes the first language.  Apply
Theorem~\ref{reg:pumping} to $a^nb^n$.  The nonempty pumped factor lies
in the first $n$ letters and is therefore $a^d$ with $1\leq d\leq n$.
Taking $t=0$ accepts $a^{n-d}b^n$, whose two counts differ, a contradiction.
For palindromes intersect a hypothetical DFA with the regular language
$a^*ba^*$.  The intersection is exactly $\{a^rba^r:r\in\N\}$.
For its own number $n$ of states, pump $a^nba^n$ at the first $n$ letters
and take $t=0$; the two $a$-blocks no longer have equal lengths.
For squares of words intersect with $a^*ba^*b$.  A square in this slice
has exactly one $b$ in each half, and each half ends in $b$, so the
intersection is exactly $\{a^rba^rb:r\in\N\}$.  Pumping its initial
$a$-block gives the same contradiction.
Finally an $n$-state unary-square DFA accepts $a^{n^2}$.  Pumping with
$t=2$ gives an accepted length $n^2+d$, where $1\leq d\leq n$.
But $n^2<n^2+d<(n+1)^2$, so this length is not a square.  Deciding square
membership itself uses the bounded search $r\leq|w|$; all four predicates
and all contradiction witnesses are primitive recursive.
\end{proof}

\subsection{Finite decision procedures and canonical minimization}

\begin{theorem}[Emptiness and shortest accepted words]\label{reg:empty}
An $n$-state DFA has a nonempty language if and only if it accepts a word
of length below $n$.  Emptiness and a shortest accepted word, when one
exists, are primitive recursive.
\end{theorem}
\begin{proof}
An accepting run can be shortened by deleting loops, preserving its
initial and final states; the remaining path has fewer than $n$ edges.
Conversely any such word witnesses nonemptiness.  Enumerate words of
length $0,1,\ldots,n-1$, lexicographically within each length, and choose
the first accepted one.  The finite enumeration tests
$\sum_{i<n}k^i$ words.  If every test fails, no accepting run exists by
the shortening argument.
\end{proof}

\begin{theorem}[Equality, inclusion, universality, and short separators]
\label{reg:equality}
For DFAs with $n,m$ states, inequality of their languages has a witnessing
word of length below $nm$.  Failure of inclusion has the same bound.
Equality, inclusion, disjointness, and universality are primitive
recursive decisions on finite automaton codes.
\end{theorem}
\begin{proof}
Use the product automaton, accepting respectively when the two bits
differ, when the first is true and the second false, or when both are
true.  Lemma~\ref{reg:run} identifies its language with the required
counterexamples.  Theorem~\ref{reg:empty} bounds every shortest witness
by $nm-1$ and decides emptiness.  Universality is emptiness of the
complement.  A negative decision produces an explicit counterexample;
a positive decision says a finite Boolean reachability table has no
accepting state.  These are complete decisions, rather than searches
that stop only on counterexamples.
\end{proof}

\begin{theorem}[Finiteness has a bounded test]\label{reg:finiteness}
An $n$-state DFA accepts arbitrarily long words if and only if it accepts
a word of length in $[n,2n)$.  If it accepts no word in this interval,
all its accepted words have length below $n$.  Thus finiteness is decidable,
and a finite accepted language can be completely enumerated.
\end{theorem}
\begin{proof}
A word of length at least $n$ has a pumping decomposition.  Repeating its
nonempty factor gives words of unbounded length; for a requested bound
$B$, choosing $t=B+1$ already makes the pumped word longer than $B$.
Conversely, given any accepted word of length at least $n$, while its
length is at least $2n$, remove the nonempty loop furnished by its first
$n+1$ run states.  The removed length is at most $n$, so the new length
is still at least $n$.  Repeat at most the original length times.  Each
successful pass decreases the length, and the result lies in $[n,2n)$.
This is a bounded primitive recursive normalization of a supplied
accepted word.  If the finite test of that interval fails, no accepted
word can have length at least $n$.  Enumerating the shorter words then
lists the whole language.  If it succeeds, the pumping construction
produces an injection of $\N$ into the language, since the lengths
$|xz|+t|y|$ strictly increase.  This gives a constructive, witnessed
meaning of infinitude without using an infinite-pigeonhole principle.
\end{proof}

\begin{proposition}[Bounded distinguishability of states]\label{reg:distinguish}
For an $n$-state DFA, whether two states have identical future acceptance
is decidable by a primitive recursive table.  If they differ, a
separating suffix of length at most $n^2$ can be computed.
\end{proposition}
\begin{proof}
Put $D_0(q,r)=F(q)\mathbin{\mathrm{xor}}F(r)$ and iterate
\[
 D_{j+1}(q,r)=D_j(q,r)\vee
   \bigvee_{a<k}D_j(\delta(q,a),\delta(r,a)).
\]
The invariant is that $D_j(q,r)=1$ exactly when a word of length at most
$j$ distinguishes $q,r$.  At a newly true entry store the responsible
letter and predecessor entry; an initially true entry stores $\eps$.
The recursion either keeps a witness or prefixes a letter to one, proving
both directions by bounded induction.
The table is monotone and has only $n^2$ bits.  If a round makes no
change, every later round is identical.  If there were a change after
round $n^2$, every earlier round would have changed as well; this would
require more than $n^2$ newly set bits.  Finite counting therefore proves
$D_{n^2}=D_{n^2+1}$.  At this fixed point, a false pair has equal accepting
bits and sends every letter to a false pair.  Induction along any supplied
word now proves identical future acceptance.  A true pair has its
recorded suffix of length at most $n^2$.
\end{proof}

\begin{theorem}[Constructive DFA minimization]\label{reg:minimize}
There is a primitive recursive procedure that takes any DFA to an
equivalent reachable DFA with the smallest possible number of states.
It returns access words for its states and distinguishing suffixes for
every pair of different states.  All these witnesses have explicit
bounds in the size of the input automaton.
\end{theorem}
\begin{proof}
First restrict to states reachable from $q_0$, using
the bounded reachability result~\ref{base:graph}; choose a shortest access word of length below $n$
for each.  Write $q\sim r$ when $D_{n^2}(q,r)=0$.
The characterization in Proposition~\ref{reg:distinguish} makes this
an equivalence relation.  More explicitly, reflexivity and symmetry are
immediate from equality of accepting bits, and transitivity follows by
comparing the bits for every word of length at most $n^2$.  Its fixed-point
equation gives $q\sim r\Rightarrow\delta(q,a)\sim\delta(r,a)$.
Choose the least reachable representative of each equivalence class by
bounded search.  These representatives, ordered numerically and
renumbered consecutively, are the states of the quotient.  Transitions
send a representative to the representative of its successor; acceptance
is its original bit.  The right-congruence property makes the choice
independent of representatives, and induction on a run proves equality
of the accepted language.
To prove minimality, let $B$ be any other DFA for that language.  For
each quotient state choose the already computed access word $u_i$ and
map it to the state reached in $B$ by $u_i$.  For $i\ne j$ the table gives
a suffix $v_{ij}$ of length at most $n^2$ accepted from exactly one of
the two original states.  If $B$ reached the same state after $u_i,u_j$,
run composition would force the same answer on $u_iv_{ij}$ and
$u_jv_{ij}$, contrary to agreement with the original automaton.  Thus
the finite map is injective.  Finite pigeonhole implies that $B$ has at
least as many states as the quotient.  Every comparison used here is a
Boolean comparison of finite runs; the universally quantified assumption
that $B$ recognizes the same language is only instantiated at the
explicitly constructed words.
\end{proof}

\begin{corollary}[Uniqueness of the reachable minimal DFA]\label{reg:unique}
Two reachable minimal complete DFAs for the same language are isomorphic
by a unique initial-state-preserving, transition-preserving map.
The map and its inverse are obtained by bounded access-word searches.
\end{corollary}
\begin{proof}
Every state of a minimal DFA is distinguishable from every other: if
not, the quotient construction of Theorem~\ref{reg:minimize} would
strictly reduce its state count while preserving the language.
For a state $q$ of the first automaton choose a shortest access word $u_q$
and map $q$ to the state reached by $u_q$ in the second.  If $u,v$ reach
the same first state, the two corresponding second states accept the
same suffixes, because both automata recognize the same language.
Distinguishability in the second automaton forces those states to be
equal.  The map therefore does not depend on the chosen access word.
It preserves the initial state, transitions, and final bits.  The symmetric
construction is an inverse, as seen by using the same access words.
Finally any map preserving the initial state and transitions must send
the state reached by $u$ to the state reached by $u$; reachability makes
this condition determine the whole map.
\end{proof}

\begin{theorem}[Effective Myhill--Nerode certificate]\label{reg:nerode}
Fix a primitive recursive language predicate $L$.  A finite right-congruence
certificate consists of a fixed PR function $c$ with values in $[m]$,
words $r_0,\ldots,r_{m-1}$ with $c(r_i)=i$, and the assertions
\[
\begin{split}
 c(u)=c(v)&\Longrightarrow c(ua)=c(va),\qquad a<k,\\
 c(u)=c(v)&\Longrightarrow L(u)=L(v).
\end{split}
\]
Such a certificate gives an $m$-state DFA for $L$.  Conversely every DFA
for $L$ supplies such a certificate after removing inaccessible states.
Consequently regularity is equivalent to the existence of this effective
finite data, with their displayed arithmetic correctness assertions.
\end{theorem}
\begin{proof}
Use initial state $c(\eps)$, transition
$\delta(i,a)=c(r_i a)$, and accepting bit $F(i)=L(r_i)$.
If the current state after $u$ is $c(u)$, then $c(r_{c(u)})=c(u)$,
so the first certificate assertion gives
$\delta(c(u),a)=c(ua)$.  Induction along the word proves the run state
is $c(w)$; the second assertion gives $F(c(w))=L(w)$.
Conversely let $c(w)$ be the state reached after $w$ in a reachable DFA,
renumbering its states below $m$ if necessary.  The run-composition lemma
proves right stability and agreement of accepting bits.  Bounded
reachability provides a representative word for every state.
This is a schema for fixed PR functions, not an assertion that all PR
programs admit a PR universal evaluator.  Nor does it infer decidable
classes or representatives from an abstract finite-index equivalence
relation on arbitrary predicates.  Those additional data are exactly
what make the converse constructive and its realization primitive
recursive.
\end{proof}

\subsection{Lean implementation obligations for this section}
\label{reg:lean}

The semantic Lean layer may use \texttt{List (Fin k)} for words and the
standard \texttt{DFA}, \texttt{NFA}, and \texttt{\ensuremath{\eps}NFA}
interfaces.  It must not be confused with the arithmetic layer.  For the
latter use a record of natural-number fields for the alphabet size, state
count, start, transition-table code, and final-bit-vector code, together
with a PR Boolean predicate \texttt{validDFA}.  The NFA code has an initial
bit-vector and a three-dimensional Boolean edge table; its epsilon variant
adds a two-dimensional edge table.  Define a total decoder for every field,
using the rejecting automaton on invalid inputs.  A regular expression is
a postfix list of tagged nodes, each child index strictly less than its
parent index. The validity predicate also requires a single final root,
exactly one parent for every nonroot node, and connectivity to that root;
thus the code describes a tree, not a shared directed acyclic graph. These
conditions are bounded PR checks.  Grammar codes are lists of the
four production tags used in Theorem~\ref{reg:grammar}.

The following are the substantive proof obligations; none is silently
assumed from an algorithm's successful execution.
\begin{enumerate}
\item \texttt{runCode\_append}: decoded iteration equals list evaluation,
with the invariant ``the first $i$ scanned letters lead to the stored
state.''  Construct the code for this equality in the object calculus.
\item \texttt{shortenPath\_correct}: verify every deletion, preserve endpoints,
and combine the explicit finite-pigeonhole witness with the decreasing
length counter.  The output bound is $n-1$.
\item \texttt{epsilonClosure\_correct} and \texttt{subset\_correct}: prove the
closure table's witness invariant, then the subset vector's prefix
invariant.  The bounded epsilon-path bound is
$\ell+(\ell+1)(n-1)$.
\item \texttt{closure\_constructors\_correct}: each closure operator must
return a valid finite code and maps in both directions between the relevant
finite path certificates.  For quotients the crucial search bound is $nm$.
\item \texttt{regexTable\_correct}, \texttt{thompson\_correct}, and
\texttt{eliminate\_correct}: use the common finite table of expression
nodes, word intervals, and state endpoints.  For elimination also record
the internal-vertex threshold and the expression-size majorant.  Structural
induction in ambient Lean alone is insufficient evidence of an \ipra{}
proof; the table invariant is the route to quantifier-free induction.
\item \texttt{regularGrammar\_correct}: implement stepwise derivation/path
conversions and prove the terminal-prefix/nonterminal-suffix invariant.
\item \texttt{pump\_correct}: extract the repeated indices from the first
$n+1$ run states and prove the loop equation by quantifier-free induction
on the pumping exponent.  State the nonregularity results as explicit
refutations of alleged DFA recognition equations.
\item \texttt{equivalence\_decision\_correct},
\texttt{finiteness\_decision\_correct}, and
\texttt{minimize\_correct}: return either bounded witnesses or checked
negative reachability tables.  For minimization add the $n^2$-round
stabilization proof, representative correctness, and the finite injection
into any competing DFA.
\item \texttt{finiteRightCongruence\_toDFA}: treat the fixed PR code $c$ and
its object-theoretic correctness derivations as parameters.  Keep the
schema separate from the internal quantification over finite DFA data.
\end{enumerate}

For example, after defining the constructors the semantic target for
subset construction is
\[
 \forall M,w\quad
 \operatorname{validENFA}(M)=1\ \longrightarrow\
 \operatorname{acceptDFA}(\operatorname{determinize}(M),w)
 =\operatorname{acceptENFA}(M,w).
\]
Both acceptance terms here are primitive recursive: the second uses the
bounded path normalization above.  The object-level target is a derivation
of this formula in the encoded \ipra\ calculus, after its Boolean and
bounded quantifiers have been compiled away.  Existing mathlib declarations
such as \texttt{NFA.toDFA\_correct} and
\texttt{\ensuremath{\eps}NFA.toNFA\_correct} are useful semantic comparison
lemmas; they do not themselves discharge that target.  The inspected
regular-expression documentation explicitly describes its syntax and
semantics; no version-independent assumption about a preexisting complete
Kleene-equivalence implementation is needed here \cite{MathlibRegex}.

\section{Finite monoids, syntactic contexts, and linear grammars}
\label{sec:monoids}

A finite monoid is a supplied multiplication table $\mu:[m]^2\to[m]$,
$m\ge1$, and an element $e<m$, satisfying
\[
 \mu(e,x)=\mu(x,e)=x,\qquad
 \mu(\mu(x,y),z)=\mu(x,\mu(y,z)).
\]
These laws and table well-formedness are decidable bounded assertions. A
recognizer additionally supplies letter images $\lambda(a)<m$ and accepting
bits $P(x)$. Its word evaluation starts at $e$ and multiplies the letter
images from left to right. All structures here are finite numerical tables;
no quotient of an arbitrary set of words is assumed.

\begin{theorem}[Finite-monoid recognition and automata]\label{mon:recognition}
A language has a supplied finite-monoid recognizer if and only if it has a
finite DFA. Both conversions are primitive recursive. A monoid with $m$
elements gives an $m$-state DFA; an $n$-state DFA gives a monoid with at most
$n^n$ elements.
\end{theorem}
\begin{proof}
From a monoid take state set $[m]$, initial state $e$, transition
$x\mathrel{\smash{\xrightarrow{a}}}\mu(x,\lambda(a))$, and final bits $P$.
Its run is exactly the monoid evaluation. Associativity and the unit laws
also prove $\widehat\lambda(uv)=\mu(\widehat\lambda(u),\widehat\lambda(v))$
by induction along $v$.

Conversely represent a transformation of $[n]$ by its $n$-entry vector.
There are $n^n$ vectors. For transformations $f,g$, define
\[
 (f\cdot g)(q)=g(f(q));
\]
thus multiplication means first $f$, then $g$. The identity vector is a
two-sided unit, and both ways of multiplying three vectors have entry
$h(g(f(q)))$. This proves the monoid laws by a bounded comparison of entries.
Send $a$ to the transition vector $q\mapsto\delta(q,a)$. Run composition
shows that the image of $w$ is
$T_w(q)=\delta^*(q,w)$ and $T_{uv}=T_u\cdot T_v=T_v\circ T_u$.
Mark vector $f$ accepting exactly when $f(q_0)$ is an accepting state of
the DFA. This gives precisely the original accepting bit. Vector encoding,
composition, and the finite multiplication table are primitive recursive.
Their proofs compare computed entries, so the required inductions are
quantifier-free after bounded Boolean compilation.
\end{proof}

\begin{theorem}[Computing the generated transition monoid]\label{mon:generated}
From an $n$-state DFA one can compute its generated transition monoid,
its multiplication table, and a representative word for each element.
It has at most $N=n^n$ elements, and every representative may be chosen
with length below $N$.
\end{theorem}
\begin{proof}
Begin with the identity vector and repeatedly add $f\cdot T_a$ for every
currently present vector $f$ and letter $a$. After $j$ rounds the marked
vectors are exactly those represented by words of length at most $j$:
append one letter in the forward direction and remove the last letter in
the reverse direction. Store a first predecessor and letter for each new
vector. The universe has $N$ elements, so after $N$ rounds the bit table
is stable by bounded saturation. A shortest path from the identity to any
marked vector in this letter-transition graph has fewer than $N$ edges,
by finite path shortening. Those paths give the claimed representatives.

The stable set contains every word transformation: induction along a
supplied word uses its closure under each letter. It is closed under
multiplication as well. If $g=T_v$ is represented by its computed word,
starting at $f$ and reading $v$ remains in the stable set and ends at
$f\cdot g$. Thus the inherited table is a monoid. Renumber the marked
vectors in numerical order to obtain its finite code. Its recognizing
bits are inherited from Theorem~\ref{mon:recognition}. Completeness is
proved by the finite-table invariant and arbitrary-word induction, not
by extrapolating from sampled words.
\end{proof}

\begin{theorem}[Effective syntactic monoid and short distinguishing contexts]
\label{mon:syntactic}
Let $M$ be a reachable minimal $n$-state DFA for $L$. Define two-sided
syntactic equivalence by
\[
 u\approx_L v\quad\longleftrightarrow\quad
 \forall x,y\;\bigl(L(xuy)\leftrightarrow L(xvy)\bigr).
\]
Then $u\approx_L v$ if and only if $T_u=T_v$ as $n$-entry vectors.
When the vectors differ, one can compute distinguishing contexts with
$|x|<n$ and $|y|\le n^2$. Consequently the syntactic monoid of a supplied
regular language is effectively finite and primitive-recursively computable.
\end{theorem}
\begin{proof}
If the vectors agree, a prefix $x$ reaches some state $q$; reading $u$ or
$v$ then gives the same state, and the common suffix $y$ gives the same
accepting bit. This proves the context equation for arbitrary supplied
$x,y$.

If the vectors differ, bounded search finds $q<n$ with
$p=T_u(q)\ne T_v(q)=r$. Reachability provides an access word $x$ for $q$
with $|x|<n$. Distinct states of a minimal DFA are distinguishable;
otherwise the quotient construction in Theorem~\ref{reg:minimize} would
reduce its number of states. The table of
Proposition~\ref{reg:distinguish} therefore provides $y$, $|y|\le n^2$,
whose accepting bits from $p,r$ differ. Hence $xuy,xvy$ distinguish the
contexts. Since vector equality is decidable, this proves both directions
constructively, including the implication from the universally quantified
context equation to vector equality.

Equivalence is a two-sided congruence because multiplication of equal
vectors on either side preserves equality. More explicitly,
$T_{aub}=T_a\cdot T_u\cdot T_b$ for arbitrary words $a,b$.
The generated transition monoid therefore represents exactly the syntactic
classes, with class multiplication represented by concatenation. Given any
input DFA, first apply the primitive recursive minimization procedure,
then Theorem~\ref{mon:generated}. This returns the actual finite multiplication
table and representatives, without forming an infinite quotient. The
unbounded context quantifiers occur only in the theorem statement;
its test is equality of finitely many entries, and its negative branch
supplies the explicit bounded contexts above.
\end{proof}

\begin{proposition}[Left-linear grammars and the danger of mixing directions]
\label{mon:leftlinear}
Finite left-linear grammars, with productions $A\to Bu$ or $A\to u$,
define exactly the regular languages. Allowing both left-linear and
right-linear productions in one grammar can define a nonregular language.
\end{proposition}
\begin{proof}
Replace each left-linear production $A\to Bu$ by $A\to u^RB$, and each
$A\to u$ by $A\to u^R$. A left-linear sentential form consists of one
nonterminal followed by a terminal suffix. Reversing it gives a terminal
prefix followed by that nonterminal, and each old production corresponds
to exactly one new production. Induction along the supplied derivation
therefore identifies the new grammar's language with the reversal of the
old one's language. The new grammar is right-linear, so
Theorem~\ref{reg:grammar} and closure under reversal give a DFA for the
old grammar. Conversely convert a DFA for $L^R$ to a right-linear grammar
and reverse this construction to obtain a left-linear grammar for $L$.
All reversals and derivation transformations are finite PR scans.

For a strict separation, use
\[
 S\to aA\mid\eps,\qquad A\to Sb.
\]
Every nonterminal derivation alternates the sentential forms
$a^jSb^j$ and $a^{j+1}Ab^j$. Choosing $S\to\eps$ terminates with
$a^jb^j$; choosing the other rule and then the forced $A$ rule increments
$j$. This bounded trace invariant proves that the language is exactly
$\{a^jb^j:j\in\N\}$, with every such word obtained by $j$ iterations
and the final empty production. It is nonregular by
Corollary~\ref{reg:nonregular}. Each production is individually linear,
but their two directions cannot be mixed while retaining the regular-language
characterization.
\end{proof}

\section{Context-free grammars and pushdown automata}
\label{sec:context-free}

This section treats a grammar as finite input data.  A variable or a terminal
is a natural number carrying a tag, a rule is a finite list, and a derivation
is an explicitly supplied finite certificate.  No set of all derivable words
is an object of arithmetic.  The notation $w\in L(G)$ abbreviates the
first-order assertion that a checked derivation certificate exists.  The
normalization and parsing theorems below subsequently replace this assertion
by one primitive recursive Boolean computation.  The conventional grammar,
pushdown, and pumping results provide the subject matter; the bounds and
certificate transformations given here supply their arithmetic content
\cite{Sipser2020}.

\subsection{Finite syntax and finite certificates}

\begin{definition}[Grammar and certificate conventions]
\label{cf:definition}
A context-free grammar is $G=(\Sigma,V,P,S)$, where
$\Sigma=[a]$, $V=[m]$, $m\geq1$, $S<m$, and $P$ is a finite list of
rules $A\to\beta$ with $A\in V$ and
$\beta\in(\Sigma\sqcup V)^*$.  Tags keep the two alphabets disjoint.
A sentential form is a word over this tagged alphabet.  A rewriting step
replaces one specified occurrence of $A$ by the right side of a specified
rule.  A derivation certificate lists the sentential forms, rule indices,
and replacement positions.  Its checker is a bounded scan of those lists.

A completed parse tree has a variable at its root.  A variable node carries
a rule index, and its ordered children have exactly that rule's right-side
labels.  Terminal nodes have no children.  A variable using an empty rule
also has no children; it contributes no letter to the yield.  The yield is
the left-to-right list of terminal leaves.  Trees can be coded as postorder
node tables: children have smaller indices, every nonroot node occurs in
exactly one child list, and the last node is the root.  These finite conditions
exclude sharing and disconnected components.  A bottom-up table computes
the yield of every node.  Write $\operatorname{Tree}(G,A,w,t)$ for the
Boolean assertion that $t$ is such a tree with root $A$ and yield $w$.
\end{definition}

Malformed data is assigned a fixed rejecting value.  Thus all operations in
this section are total functions on natural-number codes, even when their
mathematical specifications are stated only for well-formed input.

\begin{theorem}[Derivations and trees, with bounded translations]
\label{cf:tree-derivation}
There are primitive recursive translations in both directions between
derivations $A\Rightarrow^* w$ with terminal $w$ and completed parse trees
rooted at $A$ with yield $w$.  If every right side has length at most $r$, a
$k$-step derivation produces a tree with at most $1+kr$ nodes.  A tree with
$N$ nodes produces a leftmost derivation with at most $N$ steps, and all its
sentential forms have length at most $N$.
\end{theorem}

\begin{proof}
Begin with an unexpanded root labelled $A$.  Maintain the invariant that the
ordered frontier, counting unexpanded variables but ignoring expanded
empty-rule nodes, is the current sentential form.  The supplied replacement
position identifies one unexpanded variable on this frontier.  Attach the
specified right side as its ordered children.  This changes the frontier
by precisely the certified rewriting step and introduces at most $r$
nodes.  After $k$ steps the frontier is terminal, so every remaining frontier
node is a terminal and the tree is completed.  The node count follows by
adding the numbers of newly attached children.

In the other direction, visit variable nodes in preorder.  Expand the
leftmost unexpanded variable by the rule stored at that node, skipping
terminal nodes.  Its children occur before every later frontier variable,
so the next expansion again is leftmost.  Each variable is expanded once.
Every frontier consists of distinct nodes of the original tree, so its
length is at most $N$.  The final frontier is the yield.  Both procedures
use a finite work list; the respective loop bounds are $k$ and $N$.
Their frontier invariant is an equality between coded finite lists.
\end{proof}

\begin{lemma}[Nullable variables and bounded empty trees]
\label{cf:nullable}
The variables deriving the empty word are computed by $m$ rounds of
Boolean saturation.  The computation returns an explicit empty parse tree
for every variable it marks.
\end{lemma}

\begin{proof}
Put $E_0=\varnothing$ and obtain $E_{i+1}$ by adjoining every left side of a
rule whose right side contains no terminal and whose variables all belong
to $E_i$.  Keep previously marked variables.  Store the first rule that
marks each new variable and the already stored certificates for its children.
At a strict stage the number of marked variables increases.  An equality
stage is permanent, because applying the same deterministic update to the
same set gives the same result.  Consequently $E_m=E_{m+1}$: otherwise the
first $m+1$ updates would be strict in an $m$-element set.

The stored rules give empty trees by copying the previously stored child
trees.  Conversely, scan any completed empty tree in postorder.  A terminal
leaf would contribute a letter, so none occurs.  When its children have
been shown to lie in the fixed point $E_m$, the rule at a variable node
puts that node's label in $E_m$ as well.  In particular its root is marked.
If $r=\max(1,\max_{P}|\beta|)$, each stored certificate has height at most
$m+1$ and fewer than $(r+1)^{m+2}$ nodes.  This generous bound covers unary
rules and empty right sides.  Thus certificate production itself is
primitive recursive, without an unbounded search for a derivation.
\end{proof}

\begin{theorem}[Productivity, emptiness, and a bounded positive witness]
\label{cf:emptiness}
There is a primitive recursive decision procedure for $L(G)=\varnothing$.
If the language is nonempty, it produces a terminal word and a parse tree
for it of fewer than $(r+1)^{m+2}$ nodes.
\end{theorem}

\begin{proof}
Use the preceding saturation with one change: terminals in a right side
are allowed and require no membership test.  A rule becomes usable when
all of its variable occurrences are already marked.  Record a terminal
leaf for each terminal occurrence and the stored tree for each variable
occurrence.  The same increase argument gives a fixed point after $m$
rounds and the same bound gives a tree for each marked variable.  Conversely,
the postorder scan of any terminal-yield tree marks all its variable labels
in that fixed point.  The start variable is marked exactly when a derivation
certificate exists.  A negative answer therefore refutes every supplied
certificate; it does not assert failure of an unbounded search.
\end{proof}

\subsection{Normalization with witness-preserving transformations}

\begin{theorem}[Eliminating empty rules]
\label{cf:epsilon-removal}
One can transform $G$ primitively recursively into an equivalent grammar
whose only possible empty rule belongs to a fresh start variable $S_0$, and
in which $S_0$ occurs on no right side.  Both implications of language
equivalence are witnessed by primitive recursive tree transformations.
\end{theorem}

\begin{proof}
First introduce $S_0\to S$.  Compute the nullable variables.  For each rule
$A\to X_1\cdots X_k$, enumerate the $2^k$ masks of its positions.  A mask may
delete a position only when it contains a nullable variable.  Retain every
resulting nonempty right side.  Retain $S_0\to\eps$ exactly when $S$ was
nullable.  Record, with each new rule, one original rule and one deletion
mask that produced it; a bounded first-match search supplies this provenance.
No rule introduces $S_0$ on its right side.

Given an old tree with nonempty yield, process it bottom-up and delete each
child subtree with empty yield.  Every such child is a nullable variable,
so the surviving children are an enumerated new rule.  Nodes with positive
yield retain at least one child.  The root therefore yields a new tree with
the same word.  An old empty tree at the start is replaced by the recorded
rule $S_0\to\eps$.

Conversely, for each new rule use its provenance and restore every erased
child with the fixed nullable tree from Lemma~\ref{cf:nullable}.  Recurse
over the retained variable children.  The inserted trees have empty yield,
so the result has the same terminal yield.  For $S_0\to\eps$ use its stored
old empty tree.  This is a bounded traversal of the supplied tree, with a
grammar-dependent bound on the size inserted at every node.
\end{proof}

\begin{theorem}[Eliminating unit rules]
\label{cf:unit-removal}
Unit rules $A\to B$ can be removed primitively recursively, preserving
the language and the preceding convention on the start and empty rule.
\end{theorem}

\begin{proof}
For the current $M$ variables form the directed graph of unit rules.
Compute its reflexive transitive closure by paths of length at most $M-1$.
This bound is sufficient: on any longer path, find two equal vertices by
bounded search and erase the segment between them; at most the original
path length many erasures leave a path with no repeated vertex.  Record a
short path for each reachable pair.

For every reachable pair $(A,B)$ and every nonunit rule $B\to\gamma$, add
$A\to\gamma$.  In an old completed tree, contract each maximal chain of
unit nodes into one application of its last, nonunit rule.  Each such chain
must end because the input is a finite completed tree.  The resulting rule
is among the added rules.  In a new tree, expand each rule using the stored
short unit path followed by its old nonunit rule.  These operations preserve
yield.  The only empty rule before this operation belonged to $S_0$; since
no unit edge enters $S_0$, no other variable acquires an empty rule.
\end{proof}

\begin{theorem}[Effective Chomsky normal form]
\label{cf:cnf}
There is a primitive recursive function $G\mapsto C(G)$ with
$L(C(G))=L(G)$ such that every rule of $C(G)$ is
\[
 A\to BC,\qquad A\to a,\qquad\text{or, exceptionally, } S_0\to\eps,
\]
and $S_0$ never occurs on a right side.  The function returns tree
translations certifying both directions of this equality.
\end{theorem}

\begin{proof}
Apply the preceding two transformations.  In every right side of length
at least two, replace a terminal $a$ by a fresh variable $T_a$ with rule
$T_a\to a$.  Introduce at most one $T_a$ for each letter.  A long right side
$X_1\cdots X_k$, $k\geq3$, is replaced by
\[
 A\to X_1D_1,\quad D_1\to X_2D_2,\quad\ldots,\quad
 D_{k-2}\to X_{k-1}X_k,
\]
using variables private to that rule.  Expansion and contraction of these
chains, and insertion and removal of the terminal wrappers, are mutually
yield-preserving operations on the relevant tree pieces.  Unit and empty
rules are not reintroduced.  Private auxiliary variables cannot interact
with other rules, which proves the reverse implication as well.

Here is an explicit, deliberately loose size bound.  If the input has $m$
variables, $a$ letters, $p$ rules and maximum right-side length at most
$r\geq1$, put
\[
 M=m+1,\qquad K=(p+2)2^{r+1},\qquad J=M^2K.
\]
Empty-rule elimination has at most $K$ rules, unit elimination at most
$J$, and the final grammar has at most $M+a+Jr$ variables and at most
$a+Jr$ rules.  Every transformation is a finite enumeration within these
bounds.  For example, put $H=(r+1)^{M+2}$.  An original tree with $N$ nodes
can be transformed forward using at most $2(r+1)(N+1)$ nodes: addition of
the fresh start adds one, deletion and contraction reduce size, terminal
wrappers at most double it, and binary chains multiply it by at most
$r+1$.  In the reverse direction,
\[
             (M+1)(1+rH)N+H+1
\]
is a sufficient bound.  Contracting private chains and terminal wrappers
does not increase size; restoring unit paths multiplies it by at most
$M+1$; restoring deleted children inserts at most $r$ nullable trees of
size at most $H$ per variable node.  The final $H+1$ also covers the
start's empty alternative and removal of the fresh start.  Thus the
certificate translations have explicit primitive recursive size bounds.
\end{proof}

The theorem is a language-preserving normalization.  It does not claim to
preserve the number of derivations or the ambiguity of the original grammar.
For example, unit or empty cycles may produce arbitrarily many different
old trees for one word and disappear during normalization.

\subsection{Parsing and other finite decisions}

\begin{theorem}[CYK membership, reconstruction, and its arithmetic bound]
\label{cf:cyk}
Membership in a finite context-free grammar is uniformly primitive
recursive in the grammar and the input word.  For a normal-form grammar
with $M$ variables and $P$ rules, a word of length $n>0$ is decided using
at most $M(n+1)^2$ Boolean table cells and at most
$(P+1)M(n+1)^3$ elementary rule-and-split inspections.  On acceptance a tree
with $3n-1$ nodes is reconstructed, using exactly $2n-1$ rule applications.
\end{theorem}

\begin{proof}
Let $w=w_0\cdots w_{n-1}$.  For $0\leq i<j\leq n$ and a variable $A$, fill
$D(A,i,j)$ in increasing order of $j-i$.  At length one the value is the
Boolean disjunction of $A\to w_i$ over the finite rule list.  At greater
length it is
\[
 \bigvee_{A\to BC\in P}\ \bigvee_{i<k<j}
       \bigl(D(B,i,k)\wedge D(C,k,j)\bigr).
\]
Keep the least successful rule and split as a backpointer.  All referenced
intervals are shorter than the interval being filled.  The displayed cell
and inspection bounds follow simply by allowing all $M$ variables, all
$(n+1)^2$ interval endpoints, all $P$ rules, and all $n+1$ splits.

A successful terminal cell gives a one-letter tree.  A successful binary
cell joins the two trees supplied by its earlier cells.  Induction through
the table therefore constructs a valid tree with the advertised substring
as yield.  Conversely, take any valid tree for a nonempty word.  Below the
root no empty production is possible, because only $S_0$ has one and $S_0$
is absent on right sides.  At a binary node the two child yields are
nonempty.  Their boundary supplies one of the inspected splits, and a
postorder scan of the supplied tree shows that the corresponding cell is
true.  Its root supplies $D(S_0,0,n)$.

In a positive-yield normal-form tree, write $b$ for the number of binary
variable nodes and $t$ for the number of terminal-producing variable nodes.
Counting edges between variable nodes gives $2b=b+t-1$, hence $b=t-1$.
There is one terminal leaf for each terminal-producing variable, so $t=n$.
The tree thus has $(n-1)+n+n=3n-1$ nodes and $2n-1$ variable nodes.  Expanding
the backpointers with an explicit task stack constructs this many nodes;
it never searches for a tree.  For $n=0$, inspect the single permitted empty
rule.  For a general grammar, first apply Theorem~\ref{cf:cnf} and translate
the resulting certificate back.  A negative table value refutes every
supplied old certificate via the forward translation.
\end{proof}

\begin{theorem}[A direct bounded fixed-point parser for any CFG]
\label{cf:chart}
There is also a primitive recursive parser that works directly on the
original grammar, including its empty and unit cycles.  For an input of
length $n$ and $M$ variables, its chart has
$Q=M(n+1)(n+2)/2$ possible entries and stabilizes within $Q$ simultaneous
saturation rounds.
\end{theorem}

\begin{proof}
Use entries $(A,i,j)$ for $0\leq i\leq j\leq n$, allowing empty intervals.
Initially no entry is true.  At each round keep all true entries and add
$(A,i,j)$ if a rule $A\to X_1\cdots X_k$ has a nondecreasing cut vector
\[
                  i=t_0\leq t_1\leq\cdots\leq t_k=j
\]
whose pieces match its symbols.  A terminal occurrence requires
$t_\ell=t_{\ell-1}+1$ and the matching input letter.  A variable occurrence
$B$ requires the entry $(B,t_{\ell-1},t_\ell)$ from the previous round.
For an empty right side the condition is simply $i=j$.  Every rule and
cut vector lies in a finite enumeration bounded by the grammar and $n$.
An equality of two successive charts is permanent; every strict round
adds an entry.  Hence the chart after $Q$ rounds is fixed.

When an entry first appears, record its round, rule, and cut vector.
The child-variable entries used for it have strictly earlier ranks.
Recursively joining their stored trees, with terminal leaves at the
terminal positions, gives a genuine tree for the indicated substring.
The decreasing ranks prevent cycles in reconstruction even when the
grammar has cycles.  If $r\geq1$ bounds right-side lengths, fewer than
$(r+1)^{Q+2}$ nodes suffice for each reconstructed tree.

Conversely, scan any supplied parse tree in postorder, annotating each
node by its interval of terminal-leaf positions in the whole yield.
The children of a node give a nondecreasing cut vector, including equal
cuts for empty yields.  If all child-variable entries are in the final
chart, its fixed-point property includes the parent entry.  Thus every
node, including the root, is present.  The root entry $(S,0,n)$ therefore
is true exactly for generated words.  All annotations and the processed
node invariant are bounded finite data.  This proves an independent
membership algorithm without appealing to normal form.
\end{proof}

\begin{proposition}[Counting parses of a normal-form grammar]
\label{cf:parse-count}
For a normal-form grammar, the number of parse trees for a specified word
is a primitive recursive natural number.  In particular, zero, one, or
more than one parses of that word can be distinguished effectively.
\end{proposition}

\begin{proof}
Use the same interval order, replacing the Boolean disjunctions by sums and
the conjunctions by products.  At length one count the matching rule
indices; at greater length use
\[
 C(A,i,j)=\sum_{A\to BC\in P}\ \sum_{i<k<j}
                  C(B,i,k)C(C,k,j).
\]
A tree of the given interval has a unique root rule, a unique boundary
between its child yields, and an ordered pair of child trees.  Conversely
any such data joins to a tree.  The finite disjoint union and product count
is therefore exact.  Rules may either be deduplicated first or treated as
distinct indexed productions consistently with the tree convention.
There are only $3n-1$ nodes in a tree for $n>0$, which also directly gives
a finite enumeration bound.  Empty input is handled by counting the
start's empty rules.  This proposition concerns the normal-form grammar,
not the number of parses of an arbitrary grammar before normalization.
\end{proof}

\subsection{Closure constructions}

When a finite list of factor words is supplied with membership assertions
instead of factor trees, Theorem~\ref{cf:cyk} supplies the latter uniformly:
run its primitive recursive certificate constructor on every factor.
Its completeness theorem proves each resulting tree valid, and its
explicit size bound bounds the entire finite list of trees.  Thus the
tree-list constructions below also establish the ordinary membership
formulations of star and substitution.  They do not use an unbounded
finite-collection or choice principle for existential derivation witnesses.

\begin{theorem}[Union, concatenation, and star]
\label{cf:basic-closure}
Finite CFGs can be constructed primitively recursively for
$L(G_1)\cup L(G_2)$, $L(G_1)L(G_2)$, and $L(G_1)^*$.  Tree and word-factor
certificates are transformed in both directions.
\end{theorem}

\begin{proof}
Rename variable sets into disjoint finite intervals.  For union add a fresh
start $T$ and $T\to S_1\mid S_2$; its first rule selects exactly one of the
two old trees.  For concatenation use $T\to S_1S_2$.  Its two ordered subtrees
give a factorization $w=uv$, and any supplied pair of trees for $u,v$ joins
under this rule.

For star add $T\to\eps\mid S_1T$.  Repeatedly inspect the second child of a
nonempty-rule $T$ node.  Because the tree is finite, this process ends at
$T\to\eps$ after at most its node count many steps.  The first children give
a finite list of old-language factors, possibly empty, whose concatenation
is its yield.  Conversely, a supplied finite list of factor trees builds
this chain by recursion on the list length.  Empty factors cause no
termination problem: termination is measured by the supplied list or tree,
not by the length of its yield.
\end{proof}

\begin{theorem}[Reversal]
\label{cf:reversal}
Replacing the right side of every rule by its reversal gives a CFG for
$L(G)^R$.
\end{theorem}

\begin{proof}
Reverse the children at each variable node and perform the same operation
in every child subtree.  The resulting rule is a rule of the reversed
grammar.  A bottom-up calculation uses
$(x_1\cdots x_k)^R=x_k^R\cdots x_1^R$ to show that the yield is reversed.
Applying the operation twice restores the original tree.  It preserves
the node count and is implemented by one bounded traversal.
\end{proof}

\begin{theorem}[Arbitrary string homomorphisms]
\label{cf:homomorphism}
Let $h:\Sigma^*\to\Delta^*$ be specified by a finite table of words
$h(a)$, with empty images allowed.  Replacing each terminal occurrence
$a$ in a rule by the terminal block $h(a)$ gives a CFG for $h(L(G))$.
\end{theorem}

\begin{proof}
Retain each rule's old index as provenance.  In an old tree, replace each
terminal block in a variable node's rule by its image block, recursively
transforming variable children.  The new yield is the concatenation of the
letter images in the old yield, namely $h(w)$.

Conversely, inspect a new tree node and its stored old rule.  That old rule
specifies exactly which child positions arose from variable occurrences
and which consecutive terminal positions arose from each $h(a)$.  Replace
each latter block by one old terminal leaf $a$, even if that block has
length zero.  Recurse on the variable positions.  This gives an old tree
with yield $v$ and proves that the new yield is $h(v)$.  Ambiguity between
different old rules with the same image is resolved by the recorded rule
index, using no choice principle.  If $H=\max_a|h(a)|$, each original
terminal occurrence creates at most $H$ terminal nodes, so the tree-size
increase has an elementary bound in $H$ and the original size.
\end{proof}

\begin{theorem}[Context-free substitution]
\label{cf:substitution}
Suppose a CFG $H_a$ over $\Delta$ is supplied for every $a\in\Sigma$.
There is a CFG for all words $v_1\cdots v_k$ where
$a_1\cdots a_k\in L(G)$ and $v_i\in L(H_{a_i})$ for all $i$.
\end{theorem}

\begin{proof}
Rename all variable sets disjointly.  In the rules of $G$, replace a terminal
$a$ by the start variable of $H_a$, and include every rule of every $H_a$.
Substitute a supplied $H_a$ tree for every terminal leaf of a supplied $G$
tree to obtain a new tree with the desired concatenated yield.

For the reverse transformation, a variable from a component $H_a$ can
only have descendants from that component and terminals of $\Delta$;
none of its rules returns to a $G$ variable.  In a new tree, contract each
maximal component subtree entered from a $G$ rule to the old terminal $a$.
The contracted tree is an old $G$ tree, and the contracted pieces are
exactly the required factor trees.  The boundaries and factor order are
found by a finite scan of the tagged nodes.  This argument also handles
empty component yields and empty source words.
\end{proof}

\begin{theorem}[Intersection with a regular language]
\label{cf:regular-intersection}
Given a CFG $G$ and a DFA $D$ with $d$ states, one can construct a CFG for
$L(G)\cap L(D)$ primitively recursively.  If $G$ is in normal form with
$M$ variables, the product has at most $Md^2+1$ variables.
\end{theorem}

\begin{proof}
Use variables $[p,A,q]$, intended to generate words derived from $A$ that
take the DFA from $p$ to $q$.  For every binary rule add
\[
 [p,A,q]\to[p,B,r][r,C,q]\qquad(A\to BC),
\]
for all $p,r,q$.  For $A\to a$ add
$[p,A,\delta(p,a)]\to a$.  If $S_0\to\eps$ is present add
$[p,S_0,p]\to\eps$.  A fresh start $T$ has rules
$T\to[q_0,S_0,f]$ for each final state $f$.

Given an old tree and a DFA run on its yield, assign to every subtree its
entry and exit states in that run.  The boundary between the two child
yields gives $r$, so this decoration is a new parse tree.  Conversely,
erase the state tags in a new tree.  At every node the recorded endpoints
compose exactly as DFA runs compose under concatenation.  A bottom-up
scan thus verifies the intended interpretation of $[p,A,q]$; at the fresh
start it gives a run from $q_0$ to a final state.  Empty yield is handled
by the displayed diagonal rule.  All state assignments and rule generation
range over explicit finite tables.
\end{proof}

\begin{corollary}[Disjointness and a one-sided inclusion decision]
\label{cf:regular-decisions}
For a CFG $G$ and a DFA $D$, disjointness
$L(G)\cap L(D)=\varnothing$ and inclusion $L(G)\subseteq L(D)$ are primitive
recursively decidable.  Equality of a CFG language with a supplied finite
language is also primitive recursively decidable.
\end{corollary}

\begin{proof}
For disjointness apply the preceding construction and the productivity
test.  For inclusion replace the DFA's final-state set by its complement
and test emptiness of the resulting intersection.  A negative result
returns an actual counterexample word by Theorem~\ref{cf:emptiness}.
For a finite list $F$, construct its finite trie automaton, test
$L(G)\subseteq F$ in this way, and test each $w\in F$ by CYK for the other
inclusion.  Every membership predicate involved is Boolean, so these
finite tests justify the stated alternatives intuitionistically.  This
does not give a decision procedure for $L(D)\subseteq L(G)$ for an arbitrary
regular language, or for equality of two arbitrary CFG languages.
\end{proof}

\begin{theorem}[A finite EBNF-to-CFG compiler]
\label{cf:ebnf}
Consider finite grammar definitions whose right sides are finite expression
trees built from terminal letters, references to named variables, the empty
word, binary alternation and concatenation, optional expressions, and
star or positive repetition.  Interpret recursive references by finite
parse certificates.  There is a primitive recursive translation to a CFG
with the same terminal yields, with certificate transformations in both
directions.  If there are $m$ named variables, $p$ definitions, and $e$
expression-node occurrences, the result has $m+e$ variables and at most
$p+2e$ rules, all of right-side length at most two.
\end{theorem}

\begin{proof}
Give every expression-node occurrence $E$ a fresh private variable $X_E$.
For a definition $A:=E$, add $A\to X_E$.  Compile the expression nodes by
the following finite table, where $F,H$ denote their child occurrences:
\[
\begin{array}{c|l}
 E & \text{rules for }X_E\\ \hline
 \eps & X_E\to\eps\\
 a & X_E\to a\\
 \text{reference }B & X_E\to B\\
 F\mid H & X_E\to X_F\mid X_H\\
 FH & X_E\to X_FX_H\\
 F\text{ optional} & X_E\to\eps\mid X_F\\
 F^* & X_E\to\eps\mid X_FX_E\\
 F^+ & X_E\to X_F\mid X_FX_E
\end{array}
\]
Each occurrence contributes one variable and at most two rules, giving
the stated bounds.  Distinct tags keep expression occurrences, original
variables, and terminal letters separate even when their printed names
coincide.  Construction is a bounded scan of the finite syntax tree.

Here the source certificate has explicit constructors: a named variable
chooses one of its definitions; a terminal or empty expression has its
specified yield; a reference carries a certificate for the named variable;
alternation records its chosen branch; concatenation carries an ordered
pair; an optional expression records absence or one child; a star carries
a finite list of child certificates; and a positive repetition carries
a nonempty such list.  This is a finite tagged tree with finite list
fields.  Its local constructor conditions and yield equations are checked
in postorder, just as for CFG trees.  Cyclic references do not change the
definition: only an actually supplied finite certificate is accepted.

Translate a source certificate according to the table.  A star list
becomes the rightward chain of $X_E\to X_FX_E$ rules, ending in
$X_E\to\eps$.  For a nonempty repetition list, use that chain up to its
last member and end in $X_E\to X_F$.  The other constructors correspond
directly to the displayed rules.  The concatenation of the child yields
is unchanged at every node, so the resulting CFG tree has the same word.

Conversely, inspect a completed CFG tree.  Every private expression
variable has exactly the rules prescribed for its occurrence, so its
rule index identifies the original constructor.  For a star or positive
repetition, scan its $X_E$ rightward chain to obtain the finite list of
$X_F$ subtrees.  The chain ends within the supplied tree's node count;
positive repetition cannot terminate by an empty rule.  Recursively
translate these subtrees and the other constructor children.  Original
variable nodes recover the selected definition.  This constructs a valid
source certificate with the same yield.

Both traversals are bounded by their supplied certificate tables; every
source node or list cell emits a bounded number of new nodes.  If only
source-language factor membership assertions are given, the compiled
grammar's primitive recursive parser supplies the finite witnesses as
above.  Thus optionality and repetition introduce no infinitary proof
principle.  Ordinary BNF is the fragment using named variables, terminal
strings, alternatives, and explicit empty right sides; its CFG interpretation
is exactly the corresponding part of this construction.
\end{proof}

\subsection{Pushdown automata and inverse homomorphisms}

\begin{definition}[The exact pushdown convention]
\label{cf:pda-definition}
A nondeterministic PDA has a finite state set $Q$, a finite stack alphabet
$\Gamma$, an initial state $q_0$, one initial stack symbol $Z_0$, and a
finite transition list.  A transition
$(p,b,X)\mapsto(q,\beta)$ consumes $b\in\Sigma\cup\{\eps\}$, removes
the top symbol $X$, and replaces it by $\beta\in\Gamma^*$, written with
its new top on the left.  No transition is available on an empty stack.
Acceptance means: some supplied finite run consumes the entire input and
finishes with empty stack, in any state.  Epsilon transitions consume no
input.  A run checker examines a finite sequence of configurations and
transition indices; it is primitive recursive.
\end{definition}

\begin{theorem}[From CFGs to PDAs]
\label{cf:cfg-to-pda}
A one-state PDA in this convention can be constructed from a CFG, with
exactly the same language.  Finite acceptance certificates translate in
both directions primitively recursively.
\end{theorem}

\begin{proof}
The stack alphabet is the tagged union of variables and terminals, and the
initial stack is $S$.  For $A\to\beta$ use the epsilon transition replacing
top $A$ by $\beta$.  For each terminal $a$, use the transition consuming
$a$ and removing top $a$.  A preorder traversal of an old parse tree expands
its variable nodes and then matches its terminal leaves in order.  It uses
one transition per variable node or terminal leaf, and ends with empty
stack after reading exactly the yield.

In any accepting PDA run maintain the sentential form consisting of the
already consumed terminal prefix followed by the current stack.  An
epsilon expansion is exactly an old grammar rewriting step at the first
remaining variable position.  A terminal-matching step leaves this whole
sentential form unchanged, moving its split point one place to the right.
At the end it is the input word.  Deleting the terminal-matching steps
therefore gives a grammar derivation, which Theorem~\ref{cf:tree-derivation}
converts to a tree.  Both verifications are bounded scans of the given
tree or run.
\end{proof}

\begin{theorem}[From PDAs to CFGs, without a push-length assumption]
\label{cf:pda-to-cfg}
For every PDA in Definition~\ref{cf:pda-definition}, a finite CFG with the
same language and primitive recursive certificate translations can be
constructed.  If there are $q$ states, $g$ stack symbols, $t$ transitions,
and no transition pushes more than $r$ symbols, the grammar has
$q^2g+1$ variables and at most $q+t(q+1)^{r+1}$ rules.
\end{theorem}

\begin{proof}
Use variables $[p,X,s]$.  They describe a run segment that starts in state
$p$ with stack $X\gamma$, finishes in state $s$ with stack $\gamma$, and
does not expose $\gamma$ before its last configuration.  This last condition
is essential: the segment must remove this occurrence of $X$ and everything
placed above the untouched suffix.

For a transition $(p,b,X)\mapsto(u,Y_1\cdots Y_k)$ and every tuple
$s_1,\ldots,s_k\in Q$, add
\[
 [p,X,s_k]\to b\,[u,Y_1,s_1][s_1,Y_2,s_2]\cdots
                              [s_{k-1},Y_k,s_k].
\]
An epsilon $b$ contributes no terminal.  When $k=0$, instead add only
$[p,X,u]\to b$.  A fresh start has rules
$T\to[q_0,Z_0,s]$ for all $s\in Q$.  The finite rule count is bounded by
the displayed expression.

From a parse tree, execute its root transition, then the child run
segments in their order.  The first child removes $Y_1$ while leaving
$Y_2\cdots Y_k\gamma$ untouched; the next removes $Y_2$, and so forth.
Their terminal yields concatenate to exactly the consumed input of the
whole segment.  The start rule gives an accepting run.

Conversely, consider a supplied proper segment for $X$.  After its first
transition pushes $Y_1\cdots Y_k$, locate the first configuration exposing
$Y_2\cdots Y_k\gamma$, then the first exposing
$Y_3\cdots Y_k\gamma$, and so on.  These are bounded first-occurrence
searches within the supplied run.  They exist because each transition
removes only one old top symbol, so the stack height can fall by at most
one per step.  Before each first occurrence the protected suffix has not
been touched; the resulting pieces are proper segments for $Y_i$.  Their
boundary states supply precisely one enumerated grammar rule.  If $k=0$,
properness forces the segment to end immediately after its first step.
Recursing on the strictly shorter pieces constructs the tree.

This recursion can be implemented by a work list of subintervals of the
original run.  Each recursive node accounts for a different transition,
so there are no more nodes of this kind than run steps.  Terminals add at
most one leaf per transition.  With the start node, an accepting $N$-step
run gives a tree with at most $2N+1$ nodes.  Conversely the number of PDA
steps is at most the number of variable nodes of a grammar tree.  These
bounds make both transformations primitive recursive.
\end{proof}

\begin{corollary}[PDA membership with a bounded certificate constructor]
\label{cf:pda-membership}
PDA membership, including machines with epsilon cycles, is uniformly
primitive recursively decidable.  On acceptance a primitive recursive
function returns an accepting run; its length has a primitive recursive
bound in the machine and word codes.
\end{corollary}

\begin{proof}
Convert to a grammar, normalize it, and apply CYK.  Translate an accepted
normal-form tree back through the explicitly bounded constructions, then
through Theorem~\ref{cf:pda-to-cfg}.  The composition of their bounds is
primitive recursive.  A failed parsing table rules out every PDA run by
the forward certificate transformations.  The procedure never needs to
bound or traverse arbitrary epsilon computations of the original machine.
\end{proof}

\begin{theorem}[Inverse homomorphisms, including erasing maps]
\label{cf:inverse-homomorphism}
For any finite homomorphism table $h:\Sigma^*\to\Delta^*$ and CFG $G$ over
$\Delta$, a CFG for $h^{-1}(L(G))$ can be constructed primitively recursively.
\end{theorem}

\begin{proof}
Convert $G$ to a PDA $M$.  Construct a new PDA with a protected bottom marker
$\bot$ and a finite buffer in its control state.  The buffer values are
the empty word and all suffixes of the finitely many words $h(a)$; they
form a finite, explicitly enumerable set.  Initially an epsilon transition
places $M$'s initial symbol above $\bot$ and chooses its initial state with
empty buffer.

When the buffer is empty, reading an actual input letter $a$ replaces it
by $h(a)$ and leaves the simulated state and stack unchanged.  Implement
this by, for every possible top symbol including $\bot$, a transition
that replaces that symbol by itself.  A simulated epsilon transition of
$M$ is allowed regardless of the buffer.  A simulated $b$-reading transition
of $M$ is allowed exactly when the buffer begins with $b$; it consumes
that buffered letter and no actual input.  Simulated transitions are
available only above $\bot$.  When the simulated stack is empty and the
buffer is empty, an epsilon transition may pop $\bot$.  It is the only
way to empty the new stack.

During a run, the loaded images, after removing the buffered suffix,
are exactly the letters consumed by the simulated transitions.  This is
an equality of finite words checked after every step.  An accepting run
has read all original input and has empty buffer, so its simulated
transitions accept $h(w)$ in $M$.  Conversely, take an accepting run of
$M$ on $h(w)$.  Load the letters of $w$ in order whenever the buffer is
empty, interleaving the supplied simulated transitions.  Empty images
are loaded and immediately leave an empty buffer; they cause no problem
even after the simulated stack has emptied, because the bottom marker
remains.  After all input and buffered symbols have been processed, pop
the marker.  The new run has the old number of simulated steps, plus
$|w|$ loading steps and two boundary steps.  This proves both implications
by explicit run transformations.  Finally convert the new PDA to a CFG.
\end{proof}

\subsection{Pumping, finiteness, and constructive separations}

\begin{theorem}[The context-free pumping lemma with a PR witness function]
\label{cf:pumping}
Let a normal-form grammar have $M\geq1$ variables and put $p=2^M$.
From a parse tree for $w$ with $|w|\geq p$ one can primitively recursively
extract $u,v,x,y,z$ such that
\[
 w=uvxyz,\qquad |vy|>0,\qquad |vxy|\leq p,
 \qquad uv^i x y^i z\in L(G)\quad(i\in\N).
\]
The last assertion has a single primitive recursive tree constructor in
$i$.  For an arbitrary CFG the same result holds with $M$ taken from
its computed normal form, and with trees translated back to that CFG.
\end{theorem}

\begin{proof}
Consider the binary skeleton consisting only of variable nodes.  A
terminal-producing variable is a leaf.  A skeleton in which every
root-to-leaf path contains at most $h$ variable nodes has at most
$2^{h-1}$ leaves: at a binary root the two bounds add, and a leaf has one
leaf.  This calculation is proved by a bounded bottom-up scan of a given
tree, storing height and leaf count.  Since the tree for $w$ has $|w|$
such leaves, a longest path contains at least $M+1$ variable nodes.

Take the last $M+1$ variable nodes on a longest path.  Two have the same
label $A$, by finite pigeonhole; choose such a pair by bounded search,
with one a proper ancestor of the other.  The upper one's subtree has
variable height at most $M+1$.  Indeed the chosen path is longest also
inside each of its subtrees, or replacing its suffix would give a longer
whole path.  Its yield therefore has length at most $2^M=p$.

Let $x$ be the lower subtree's yield, $v$ and $y$ the terminal yields to
its left and right within the upper subtree, and $u,z$ the yields outside
the upper subtree.  These are consecutive intervals of $w$, so
$w=uvxyz$ and $|vxy|\leq p$.  Between the repeated variables there is a
binary node and hence a sibling subtree.  Every such sibling has nonempty
yield, because an empty rule cannot occur below the fresh start.  Thus
$|vy|>0$.

The piece between the two $A$ occurrences is a context with one hole
labelled $A$, producing $v$ before that hole and $y$ after it.  Delete
this context for $i=0$, use it once for $i=1$, and wrap $i$ copies around
the lower tree in general.  A direct concatenation calculation gives
$v^i x y^i$ as the new upper yield.  Each splice joins equal variable
labels, so it remains a valid parse tree.  A primitive recursion on $i$
copies a fixed finite context; if the original tree has $N$ nodes,
$(i+1)N$ is a sufficient size bound.  Its validity and yield equations
are quantifier-free induction invariants.  This constructs certificates
for every $i$, rather than inferring their existence from an infinite
set argument.
\end{proof}

\begin{theorem}[Deciding finiteness and returning either kind of witness]
\label{cf:finiteness}
Finiteness of a CFG language is primitive recursively decidable.  The
procedure returns either an exhaustive finite list of its words, or a
word with a pumping decomposition that yields accepted words of arbitrarily
large length.  If $p=2^M$ for the normalized grammar, it suffices to test
all words with
\[
                  p\leq |w|<2p.
\]
\end{theorem}

\begin{proof}
If this finite test finds an accepted word, use
Theorem~\ref{cf:pumping}.  Since $|vy|>0$, the constructed word for index
$i$ has length $|uxz|+i|vy|\geq i$.  Taking $i=N+1$ supplies an accepted
word longer than any supplied bound $N$.

Suppose the test finds none.  To refute any alleged accepted tree of length
$n\geq p$, transform it to normal form.  If its length is at least $2p$,
pump down once.  The length decreases by an integer between $1$ and $p$,
so it remains at least $p$.  Repeat this bounded process at most the
initial $n$ times, leaving the tree unchanged once its length is less
than $2p$.  The decreasing-length invariant forces arrival at an accepted
tree of length in $[p,2p)$, contradicting the finite test.  Thus every
accepted word has length less than $p$.  Enumerate all such words and
retain precisely those accepted by CYK; this is the required exhaustive
list.  All word enumerations have a bound supplied by the finite alphabet
and $2p$, and the pumping-down argument uses the supplied tree's length
as its loop bound.  No least member of an unbounded set is invoked.
\end{proof}

\begin{proposition}[The equal-block language]
\label{cf:equal-blocks}
Over distinct letters $a,b$, the grammar $S\to aSb\mid\eps$ generates
exactly $E_{ab}=\{a^n b^n:n\in\N\}$.  Thus, together with its regular
pumping obstruction, it witnesses the strict inclusion of regular
languages in context-free languages.
\end{proposition}

\begin{proof}
Applying the first rule $n$ times and the second rule once is a certificate
for $a^n b^n$.  Conversely, in any completed tree each nonempty rule has
exactly one variable child.  Follow this chain until its final empty rule;
its finite number $n$ of nonempty rules supplies exactly $n$ initial $a$'s
and $n$ final $b$'s.  The traversal is bounded by the given tree size.
For completeness, if a $d$-state DFA accepted precisely this language,
the first $d+1$ states while reading the prefix $a^d$ would repeat.
Deleting the nonempty loop from the accepted word $a^d b^d$ preserves
acceptance but changes the number of $a$'s, a contradiction.  The loop is
found by finite pigeonhole.  A DFA has a nonempty state set, so $d\geq1$.
\end{proof}

The balanced-parenthesis grammar $S\to\eps\mid(S)S$ gives another
important example.  Its exact balance characterization, forced first-return
split, unique parse tree, and enumeration are proved together in
Theorem~\ref{count:dyck}, avoiding separate parsing and counting assumptions.

\begin{theorem}[An explicit separating word against every CFG]
\label{cf:abc-separator}
Let $T=\{a^n b^n c^n:n\in\N\}$ over three distinct letters.  There is a
primitive recursive algorithm which, on every finite CFG $G$ over this
alphabet, returns a word where $L(G)$ and $T$ differ.  With $M$ variables
in $C(G)$, the word has length at most $3\cdot2^M$.  In particular $T$
is not context-free, in the constructive sense that every proposed
grammar is effectively refuted.
\end{theorem}

\begin{proof}
Set $p=2^M$ and run the parser on $w=a^p b^p c^p$.  If it rejects, $w$ is
the required separating word, with its negative parsing table as
certificate.  If it accepts, extract a tree and the pumping decomposition
$w=uvxyz$.  The interval $vxy$ has length at most $p$.  An interval
containing both an $a$ and a $c$ would contain all $p$ intervening $b$'s
and at least those two additional letters, so it would have length at
least $p+2$.  Therefore $v$ and $y$ together affect at most two of the
three letter types.  They remove at least one letter because $|vy|>0$.

The pumped-down word $uxz$ is still accepted.  One letter type retains
its original count $p$, while another has smaller count.  Consequently
the three counts are unequal, so $uxz\notin T$.  Its length is at most
$3p$, and its acceptance certificate is supplied by pumping.  These two
branches are selected by one Boolean computation, not by excluded
middle for an arbitrary language predicate.  The membership and
nonmembership assertions for $T$ are just finite order and count checks.
\end{proof}

\begin{theorem}[Failure of intersection and complement closure]
\label{cf:nonclosures}
There are explicit context-free languages whose intersection is not
context-free, and an explicit context-free language whose complement
relative to $\{a,b,c\}^*$ is not context-free.
\end{theorem}

\begin{proof}
Let
\[
 L_1=\{a^n b^n c^k:n,k\in\N\},\qquad
 L_2=\{a^k b^n c^n:n,k\in\N\}.
\]
For $L_1$ use start $U\to SC$, with
$S\to aSb\mid\eps$ and $C\to cC\mid\eps$.
For $L_2$ use start $U\to AR$, with
$A\to aA\mid\eps$ and $R\to bRc\mid\eps$.
The chain argument of Proposition~\ref{cf:equal-blocks} verifies their
languages, and a word belongs to both exactly when all three block
lengths agree.  Their intersection is $T$, excluded by
Theorem~\ref{cf:abc-separator}.

One can exhibit the complement counterexample directly, avoiding any
appeal to a general classical De Morgan principle.  Write $E_{bc}$ for
the equal-block grammar over $b,c$, and let
\begin{align*}
 R_0&=\{a,b,c\}^*\setminus a^*b^*c^*,\\
 K&=R_0\ \cup\ (a^+E_{ab}\cup E_{ab}b^+)c^*
       \ \cup\ a^*(b^+E_{bc}\cup E_{bc}c^+).
\end{align*}
The language $R_0$ is regular; all the other displayed pieces are
context-free by the already proved union and concatenation constructions.
They form an explicit CFG expression, and those constructions compile
it to a finite grammar.  A word outside block order belongs to $R_0$.
A word in block order is $a^i b^j c^k$; its counts are not all equal
exactly when $i>j$, $i<j$, $j>k$, or $j<k$.  These four decidable cases
are exactly the other four pieces of $K$.  Hence $K$ is precisely the
complement of $T$, and its complement is $T$.  All equivalences here
are verified using Boolean order and count comparisons.
\end{proof}

\begin{proposition}[A bounded tape recognizer and the hierarchy separation]
\label{cf:abc-noncontracting}
The language $T$ has a deterministic recognizer using only the input's
cells, fixed endmarkers, and finitely many marked variants of its letters.
The computation has a primitive recursive quadratic time bound.  Moreover,
context-free languages are a proper subclass of the languages generated
by noncontracting grammars with the isolated-start empty-word convention.
\end{proposition}

\begin{proof}
First scan to check block order.  Repeatedly mark the leftmost unmarked
$a$, then the leftmost unmarked $b$, then the leftmost unmarked $c$,
rejecting if either latter search fails.  When no unmarked $a$ remains,
accept exactly if no unmarked $b$ or $c$ remains.  Each successful round
marks one of each letter.  At most the input length $n$ many rounds are
possible, and each search and return to the left end is a bounded tape
scan.  Fewer than ten scans per round, and finitely many initial and
final scans, give the loose move bound $20(n+2)^2$.  Permanent endmarkers
and three marked variants of the letters suffice; no additional unbounded
tape is used.  Empty input is accepted immediately.

Counting marked letters proves soundness.  If the three block counts
are equal, after each successful round their unmarked remainders are
still equal, so the next indicated searches succeed until all are zero.
This proves completeness by a bounded invariant on the rounds.

For the language-class inclusion, the normal form from
Theorem~\ref{cf:cnf} is itself noncontracting apart from its isolated
start's empty rule.  The explicit noncontracting grammar for $T$ is
proved in Proposition~\ref{eff:abc}.  Theorem~\ref{cf:abc-separator}
excludes every CFG for that same language, giving strictness.
\end{proof}

\subsection{What the arithmetic and Lean claims mean here}
\label{cf:audit}

Every induction in this section can be organized around one of four
finite records: a processed prefix of a node table, a processed prefix
of a supplied transition trace, a bounded grammar-enumeration table, or
an already filled portion of a parsing table.  Each record's validity,
list equality, label agreement, size bound, and Boolean cell value are
primitive recursive predicates.  A conjunction over all already processed
indices is itself a bounded Boolean fold.  Quantifier-free induction on
that fold proves the invariant.  Apparent tree induction in the prose is
implemented by a postorder node loop; recursive certificate construction
is implemented by its explicitly bounded work list.  The pumping index
uses ordinary primitive recursion with a fixed finite context.

In particular, the saturation arguments do not use induction on the
unbounded statement ``$A$ has some derivation.''  The positive direction
checks a stored finite witness; the negative direction transforms or
scans an arbitrary witness supplied as a parameter.  The same distinction
justifies the negative answers of CYK and the emptiness and finiteness
procedures.  All infinite-language conclusions are expressed by a uniform
word-and-certificate constructor for each requested numerical bound.

The accompanying \texttt{lean/ContextFreeOutline.lean.txt} specifies the
finite records, Boolean predicates, algorithms, theorem interfaces, and
the order of the proof obligations.  It is an outline, not a claimed
compiled formalization.  Existing Lean grammar definitions can serve as
an external semantic comparison, but using a theorem about unrestricted
Lean predicates alone would not establish its derivability in \ipra.
The object-level route is to implement the finite functions in the
repository's primitive-recursive coding layer and prove the Boolean
invariants just described.

\section{Effective languages, finite rewriting, and diagonal limits}
\label{sec:effective}

This section treats an unbounded search as an existence of a finite certificate,
not as a total function.  Its basic machine constructions belong to the tradition
of finite machine coding and universal simulation \cite{Turing1936}; all bounds
and diagonal arguments used below are supplied here.  In particular, a bounded
universal machine simulator must not be confused with a total universal
evaluator for primitive recursive functions.

\subsection{A machine model with a primitive recursive bounded interpreter}

Fix a finite tape alphabet $\{0,\ldots,b-1\}$, with $b\geq3$ and blank $0$.
A deterministic machine has nonhalting states $0,\ldots,q-1$, a halting state
$q$, and a finite table indexed by a state and a scanned symbol.  A table entry
specifies a new state, a written symbol, and a move in $\{\mathrm L,
\mathrm R,\mathrm S\}$.  Invalid codes are interpreted as a machine that
immediately halts; this convention makes all decoding operations total.

Represent a configuration by $(p,L,R)$, where the least significant base-$b$
digit of $R$ is the scanned cell and that of $L$ is the cell immediately to its
left.  Unrepresented digits are blank.  If the current instruction writes $d$,
put $T=\lfloor R/b\rfloor$.  The three tape updates are
\[
\begin{array}{c|cc}
 &L'&R'\\ \hline
\mathrm R&d+bL&T\\
\mathrm L&\lfloor L/b\rfloor&(L\bmod b)+bd+b^2T\\
\mathrm S&L&d+bT.
\end{array}
\]
The state becomes the entry's new state.  A halted configuration is unchanged.
Natural-number inputs have a fixed delimiter-free binary representation using
tape symbols $1,2$, followed by blanks; zero has the one-symbol representation
$1$.  Finite-word inputs are encoded as natural numbers by the word coding of
the foundations chapter.  The choice of this input convention has no effect
on any theorem, but is fixed throughout.

\begin{lemma}[Bounded simulation and finite traces]\label{eff:simulation}
There are uniform primitive recursive functions
$\operatorname{Step}(e,C)$ and $\operatorname{Run}(e,x,t)$, and Boolean
primitive recursive predicates
$\operatorname{Halt}(e,x,t)$ and $\operatorname{Out}(e,x,t,j)$,
expressing respectively that the machine coded by $e$ has halted by stage $t$
and has halted with Boolean output $j$.  A halting or output witness can be
converted primitive recursively to a locally checked configuration trace,
and conversely a valid trace of length $t+1$ proves the corresponding
bounded-simulation assertion.
\end{lemma}
\begin{proof}
Table lookup, tuple decoding, remainder, quotient, and the displayed updates
are primitive recursive.  Define $\operatorname{Run}(e,x,0)$ to be the initial
configuration and iterate $\operatorname{Step}$ exactly $t$ times.  Output is
read from the scanned cell at halting, with any non-Boolean value treated
as an invalid Boolean output.  Store the successive configurations during
this iteration to obtain a trace.  Its checker verifies the initial entry,
each adjacent pair, and the final state by a bounded conjunction.
For the converse, use induction on the number of checked adjacent pairs:
the computed state and the supplied trace state coincide.  The induction
invariant is an equality of codes together with a bounded Boolean check.
Thus this is quantifier-free induction after characteristic functions are
introduced.  Determinism also proves that two valid halting traces of the
same computation have the same output.  Absorption of the halting state
proves monotonicity in $t$.
\end{proof}

\begin{lemma}[Compilation of each fixed primitive recursive predicate]
\label{eff:compile}
For each fixed primitive recursive Boolean predicate $R(\vec x)$ there are
a fixed machine $m_R$ and a primitive recursive clock $c_R(\vec x)$ such that
\ipra\ proves that $m_R$ halts by that clock and outputs $R(\vec x)$.
This is a metatheoretic construction on the defining expression for $R$,
not a uniform primitive recursive evaluator for all such expressions.
\end{lemma}
\begin{proof}
First compile into finite register programs with increment, conditional
decrement, and jumps.  Zero, successor, and projections have direct programs.
Composition uses disjoint scratch registers and copying loops; a source is
copied by decrementing it while incrementing both the target and a temporary
register, then decrementing the temporary to restore the source.  Each such loop has
its supplied natural-number length as a clock.  For primitive recursion
$F(\vec x,0)=g(\vec x)$ and
$F(\vec x,n+1)=h(\vec x,n,F(\vec x,n))$, keep a counter, the original inputs,
and the current value in separate registers.  Repeat the compiled $h$
program exactly $n$ times.  If clocks for $g,h$ have already been supplied,
the clock for this loop is given by primitive recursion, adding at stage $i$
the clock for $h(\vec x,i,F(\vec x,i))$ and the explicitly bounded copying
cost.  The invariant identifies the counter and stored value with $i$ and
$F(\vec x,i)$; it is quantifier-free.

For completeness, these register programs have the claimed tape-machine
implementation.  Encode their finitely many registers by unary blocks
separated by a distinct delimiter.  A fixed finite control selects a block,
checks emptiness, and inserts or removes one mark by a finite scan and shift
of the remaining blocks.  If their total length is $s$, each instruction is
implemented in at most $C(s+1)^2$ tape steps, for a constant depending only on
the fixed register program.  In $t$ register instructions the total number
of marks grows by at most $t$.  Consequently
$C(t+1)(s+t+1)^2$, enlarged for fixed delimiters, bounds the simulation.
Converting the fixed binary input encoding to unary register blocks and
back is a primitive recursive bounded process (repeated doubling while
reading the input suffices).  Compose these clocks with the previously
constructed register clock.  The simulation invariant is equality of the
encoded register contents at instruction boundaries.  This supplies both
an actual compiler and its \ipra\ verification schema.
\end{proof}

\subsection{Recognition means a finite positive certificate}

\begin{definition}\label{eff:def-re}
A \emph{certificate presentation} of a language is a fixed primitive
recursive Boolean verifier $R(w,c)$; its membership formula is
\[
 L_R(w)\quad\Longleftrightarrow\quad \exists c\;R(w,c)=1.
\]
A \emph{primitive recursive enumeration with possible silence} is a
primitive recursive function $E(s)=(a_s,w_s)$, where $a_s\in\{0,1\}$ and
$a_s=0$ means that stage $s$ outputs nothing.  A machine recognizes a
language if membership is equivalent to existence of a finite accepting
trace.  A rejecting halt, when present, is not an accepting trace.
\end{definition}

\begin{theorem}[Three equivalent effective presentations]
\label{eff:three-presentations}
Certificate presentations, primitive recursive enumerations with possible
silence, and machine recognition admit explicit conversions in both
directions.  Each equivalence is proved in \ipra\ for the given
presentation.  The conversions on positive witnesses are primitive
recursive.
\end{theorem}
\begin{proof}
For $R$, inspect the fixed primitive recursive pair coding at each stage.
Emit silence on gaps in a sparse coding or on invalid word codes; on a valid
pair $(w,c)$, output $w$ precisely when $R(w,c)=1$.  A certificate $c$ gives the
particular stage $\langle w,c\rangle$, and an output stage supplies the pair
that was checked.  Conversely, the stage number itself is a certificate
that an enumerator outputs $w$.

A recognizer on input $w$ tests $R(w,0),R(w,1),\ldots$ in order and accepts
at the first success, using the machine of Lemma~\ref{eff:compile}.  Given
a successful $c$, its accepting trace has a primitive recursive bound:
sum the fixed clocks for the tests with indices at most $c$, and add the
counter and copying costs.  Conversely an accepting trace records which
test succeeded, so its index is recoverable by a bounded scan of that
trace.  Finally, for any machine, the predicate saying that its bounded
run has accepted is a primitive recursive verifier by
Lemma~\ref{eff:simulation}.  None of these steps asserts that an
unsuccessful search terminates.
\end{proof}

\begin{proposition}[Binary positive operations and reversal]
\label{eff:binary-closure}
Certificate-presented languages are closed under union, intersection,
concatenation, and reversal, with primitive recursive verifiers and
primitive recursive translations of the displayed witnesses.
\end{proposition}
\begin{proof}
For union a certificate contains a tag and a certificate for the indicated
operand.  For intersection it contains both certificates.  For
concatenation it contains a split point $i\leq |w|$ and certificates for
$\operatorname{take}_i(w)$ and $\operatorname{drop}_i(w)$.  A successful
concatenation witness determines these two words because concatenation is
cancellative and lengths add.  Conversely two successful witnesses and the
word equation supply that split record.  Reversal uses the verifier
$R(\operatorname{rev}(w),c)$ and involutivity of reversal.  Every check is
bounded by the supplied record.  The proofs use only elimination and
introduction of finitely many existential witnesses; they require no
principle of choice or unbounded collection.
\end{proof}

\begin{proposition}[Homomorphic images and preimages]
\label{eff:hom-closure}
For a finite-table word homomorphism $h$, both $h(L_R)$ and
$h^{-1}(L_R)$ have certificate presentations, including when $h$ erases
letters.  If $L$ has a primitive recursive characteristic function and
$h$ is nonerasing, then $h(L)$ also has one.  Every inverse homomorphic
image of a primitive recursive language is primitive recursive.
\end{proposition}
\begin{proof}
For the image, certify a preimage word $u$, a certificate $c$, the equality
$h(u)=w$, and $R(u,c)=1$.  For the preimage use $R(h(w),c)$.
For a nonerasing $h$, the equality $h(u)=w$ implies $|u|\leq |w|$ by adding
the positive lengths of letter images.  Enumerate the finitely many words
$u$ over the source alphabet with length at most $|w|$, and test both the
equality and the given characteristic function.  Their number is
$\sum_{i=0}^{|w|}a^i$ for source-alphabet size $a$; the empty alphabet is
handled by its sole word $\eps$.  This bounded disjunction is primitive
recursive.  For an inverse image simply evaluate the characteristic
function at $h(w)$.
\end{proof}

\paragraph{An essential finite-collection distinction.}
For unbounded certificates, the formula
$\forall i<k\;\exists c\,R(v_i,c)$ is not silently replaced by the
existence of one coded list of all the $c$'s.  That replacement is a
collection principle, not a consequence of the observation that each
individual certificate is finite.  The next proposition consequently
states its constructive input data explicitly.

\begin{proposition}[Certified iteration and certified substitution]
\label{eff:certified-star}
The following operations preserve certificate presentations:
\begin{enumerate}
\item Kleene iteration with a certificate consisting of a list of factors,
a list of their certificates, and the concatenation equation;
\item substitution from a certificate-presented source language and a
finite family of certificate-presented replacement languages, with a
source witness, replacement factors, and one coded list of all replacement
certificates.
\end{enumerate}
For primitive recursively decided operand membership these agree, internally
in \ipra, with
the usual bounded-cut definitions.  For general certificate languages
they agree with the usual operations on standard languages externally;
internally a common bound on all required certificates suffices to prove
the corresponding equivalence with a bounded family of existential
membership assertions.
\end{proposition}
\begin{proof}
For iteration, the verifier checks the list lengths, each $R(v_i,c_i)$,
and $w=v_0\cdots v_{k-1}$.  Empty factors can be removed together with their
certificates; for $w\ne\eps$ the remaining $k$ is at most $|w|$, while
$\eps$ is certified by the empty list.  For substitution a certificate
contains $u=a_0\cdots a_{k-1}$, a certificate for source membership,
factors $v_0,\ldots,v_{k-1}$, their verifier certificates, and
$w=v_0\cdots v_{k-1}$.  All these checks are bounded by the supplied lists.
Witness insertion, removal, and concatenation are primitive recursive.

For a decidable predicate, no unbounded certificate choices are needed:
evaluate each factor directly.  If a common certificate bound $B$ is
provided instead, find the least successful $c_i\leq B$ by bounded search
and store the resulting list by recursion on $i$.  The invariant is the
bounded conjunction of successful checks for the stored prefix.  For
standard finite lists, ordinary metatheoretic finite choice yields the
usual extensional interpretation.  We make no claim that this last
unbounded collection is an additional theorem of \ipra.
\end{proof}

\subsection{Decisions require a termination witness}

\begin{definition}\label{eff:def-total}
A Boolean machine has \emph{totality data} when its stipulated premise is
\[
 \forall w\;\exists t\;\bigl(
 \operatorname{Out}(d,w,t,0)=1\ \vee\
 \operatorname{Out}(d,w,t,1)=1\bigr).
\]
A primitive recursive clock is stronger data: a fixed primitive recursive
$b(w)$ for which the corresponding disjunction holds at $t=b(w)$.
A language is \emph{decided} by these data when output $1$ is equivalent
to membership.  The totality premise may be used as a hypothesis; the
chapter does not claim that \ipra\ proves it for every externally total
recursive machine.
\end{definition}

\begin{proposition}[What total decision supplies]\label{eff:total-decision}
A correct Boolean decider with the displayed totality data proves
$L(w)\vee\neg L(w)$ for each $w$, and gives certificate presentations for
both $L$ and $\neg L$.  A primitive recursive clock gives a primitive
recursive characteristic function.  Conversely, a primitive recursive
characteristic function supplies a clocked decider.
\end{proposition}
\begin{proof}
Eliminate the totality witness and inspect the output bit.  Output $1$
gives membership by correctness.  If the bit is $0$, a hypothetical
membership assertion would give an output-$1$ trace; deterministic output
uniqueness contradicts the existing output-$0$ trace.  This yields the
disjunction constructively.  Output traces are the certificates.  At a
supplied clock, inspect the bounded run to obtain the characteristic bit;
Lemma~\ref{eff:compile} proves the converse.  No search for a clock is
performed.
\end{proof}

\begin{proposition}[The constructive two-recognizer theorem]
\label{eff:two-recognizers}
Suppose $L$ and $\neg L$ have certificate verifiers, and suppose additionally
that $\forall w\,(L(w)\vee\neg L(w))$ is supplied.  There is a correct total
Boolean decider for $L$.  Without an explicit primitive recursive witness
bound, this conclusion does not assert a primitive recursive clock.
\end{proposition}
\begin{proof}
At stage $s$, test all certificates at most $s$ for both verifiers, stopping
with the corresponding bit at the first success.  The two sides cannot
both succeed, since their correctness would give $L(w)\wedge\neg L(w)$.
Use the supplied disjunction for the particular input.  In its left case
eliminate a positive certificate, and in its right case eliminate a
negative one.  Either yields a finite stage at which the search stops.
The primitive recursive test clocks give a finite machine trace up to that
stage, just as in Theorem~\ref{eff:three-presentations}.  This proves the
required totality implication without excluded middle.  Dropping the
exhaustiveness premise would replace this proof by an unjustified
classical step; we do not do so.
\end{proof}

\subsection{Noncontracting grammars are decided by a finite graph}

A finite rewriting grammar has disjoint finite terminal and nonterminal
alphabets, a start nonterminal $S$, and finitely many rules $u\to v$, where
$u$ is nonempty and contains a nonterminal.  One step replaces one occurrence
of $u$ by $v$ in a context.  It is \emph{noncontracting} if $|u|\leq |v|$
for every rule, with the permitted exception $S\to\eps$, provided $S$
never occurs on a right-hand side.  This is the noncontracting grammar
presentation of the context-sensitive level used here.  We do not use a
separate grammar--linear-bounded-automaton equivalence as an unstated lemma.

\begin{proposition}[Every finite grammar has recognizable membership]
\label{eff:grammar-recognizable}
For every finite rewriting grammar, with or without length-decreasing rules,
the generated language has a primitive recursive certificate verifier and
therefore a recognizer.  A supplied finite derivation is a certificate;
no bound on its length or on its intermediate word lengths is required.
\end{proposition}
\begin{proof}
A certificate is a list of sentential forms together with a rule index and
rewrite position for each adjacent pair.  Check that it starts with $S$,
ends with the proposed terminal word, and that every listed replacement
has the specified left and right contexts.  Each word and every rule index
is supplied in the record, so all checks are bounded by that record's
finite lengths.  Encoding these loops yields a primitive recursive Boolean
verifier uniformly in the grammar and certificate.  A successful check is
exactly the finite sequence required by the definition of derivability.
Conversely an actual finite derivation supplies those forms, rules, and
positions.  Apply Theorem~\ref{eff:three-presentations}.  This proves the
grammar-to-recognizer direction without asserting its more involved converse.
\end{proof}

\paragraph{The finite graph bound.}
We use the finite reachability theorem of the foundations
(Lemma~\ref{base:graph}): a graph on $V\geq1$ numbered vertices has a
primitive recursive reachability decision procedure, and a supplied path
can be shortened to fewer than $V$ edges by deleting loops.  Its algorithm
retains a simple prefix while processing the supplied path and truncates
at a repeated vertex.  The proof, including primitive recursive witness
reconstruction, has already been given there.

\begin{theorem}[Uniform primitive recursive membership for noncontracting
grammars]\label{eff:noncontracting}
There is a uniform primitive recursive decision function $D(G,w)$ for
membership in a well-formed noncontracting grammar.  For a nonempty target
of length $n$ over a total grammar alphabet of size $m$, a positive result
has a derivation with fewer than
\[
 V(G,n)=\sum_{i=0}^{n}m^i
\]
steps.  Both the Boolean answer and a positive derivation certificate can
be computed primitive recursively from $G,w$.
\end{theorem}
\begin{proof}
For $n>0$, every sentential form in a successful derivation has length at
most $n$: each ordinary rule preserves or increases length.  The exceptional
start-to-empty rule cannot occur on such a derivation, because the empty
word has no occurrence of a left-hand side and cannot subsequently become
nonempty.  Enumerate all words of length at most $n$ over the total
alphabet.  There are exactly $V(G,n)$ of them, and both the start word and
the target are among them.  Place an edge from $u$ to $v$ when some rule,
applied at a position at most $n$, transforms $u$ into $v$.  Finite rule
lookup and the position scan make this an effective finite edge table.

Apply Lemma~\ref{base:graph}.  A graph path is precisely a grammar
derivation remaining within the length bound; every successful derivation
satisfies that bound.  This proves both directions and the certificate
bound.  For $n=0$, acceptance is exactly the presence of $S\to\eps$.
Indeed no noncontracting rule reaches length zero, and $S$ cannot reappear
after it is first removed because it is absent from every right-hand side.
Reject malformed grammars by a preliminary finite check.  All enumeration
bounds, code bounds for the finite lists, and witness reconstructions are
primitive recursive uniformly in $G,w$.
\end{proof}

\begin{theorem}[An explicit primitive recursive language beyond all
noncontracting grammars]\label{eff:nc-pr-separation}
Over the binary alphabet there is a primitive recursive language $J$ and
a primitive recursive separating-word function such that, for every
noncontracting grammar code $e$, the function returns a word on which $J$
and that grammar disagree.  Consequently noncontracting grammar languages
form a proper subclass of primitive recursively decided languages.
\end{theorem}
\begin{proof}
Make the numerical convention in Theorem~\ref{eff:noncontracting} explicit.
A grammar code is a tuple containing its terminal count, nonterminal count,
start index, and finite list of pairs of word codes.  Decode the tuple and
lists by the foundations' primitive recursive operations, and perform the
finite well-formedness and noncontraction tests.  For this theorem, also
require exactly the two terminal symbols $0,1$.  On malformed codes, codes
with another terminal alphabet, and invalid target words, assign acceptance
bit $0$.  On a valid binary grammar use the decision procedure already
proved.  This gives a single total primitive recursive Boolean function
$D_2(e,w)$, and every binary noncontracting grammar has a valid code for
which that function is its characteristic function.

Let $\operatorname{bin}(e)$ be the canonical nonempty binary representation
of $e$, with zero represented by $0$.  Canonicality, decoding, and this
encoding are primitive recursive, and decoding after encoding returns $e$.
Define
\[
 J(w)=
 \begin{cases}
  1-D_2(e,w),&w=\operatorname{bin}(e)\text{ is canonical},\\
  0,&\text{otherwise}.
 \end{cases}
\]
This is a primitive recursive characteristic function.  For every natural
number $e$, substitution of its explicitly computed binary word yields
\[
 J(\operatorname{bin}(e))
     =1-D_2(e,\operatorname{bin}(e))
     \ne D_2(e,\operatorname{bin}(e)).
\]
The inequality is simply the checked two-case fact that a Boolean bit
differs from its complement.  For a valid binary grammar code it is an
explicit separator from that grammar; no search is involved.  Theorem~
\ref{eff:noncontracting} gives the inclusion of the whole class into the
primitive recursive languages, and this one $J$ proves strictness.

This diagonalization is legitimate in \ipra: the decision procedure is
uniform only over the restricted finite syntax of noncontracting grammars.
The resulting $J$ leaves that restricted class.  No universal evaluator for
all primitive recursive functions is used, and no assumption that an
externally total recursive function has a primitive recursive clock is
made.
\end{proof}

\begin{proposition}[Union, reversal, and regular intersection at the
noncontracting level]\label{eff:nc-closure}
Noncontracting grammar languages are closed constructively under union,
reversal, and intersection with a language given by a deterministic finite
automaton.  The grammar transformations and all translations of finite
derivation witnesses are primitive recursive.
\end{proposition}
\begin{proof}
For union, rename the two nonterminal alphabets to be disjoint, remove their
exceptional empty rules, and introduce a fresh start with alternatives to
the two old starts.  Add an empty rule at the fresh start exactly when one
old grammar has its exceptional empty rule.  Every nonempty derivation
stays in its chosen component, since each rule of the other component
requires a nonterminal that never occurs.  The old starts cannot reappear
inside their component, so removing their empty rules loses no nonempty
word.  For reversal replace every rule $u\to v$ by
$\operatorname{rev}(u)\to\operatorname{rev}(v)$.  Reverse each sentential
form in a trace; a context $xuy$ becomes
$\operatorname{rev}(y)\operatorname{rev}(u)\operatorname{rev}(x)$.
Lengths and the exceptional-start condition are preserved.

Here is a full regular-intersection construction, which avoids invoking an
unproved machine-equivalence theorem.  First handle the empty word
separately and delete the exceptional empty rule.  Let the automaton have
finite state set $Q$, start $q_0$, transition $\delta$, and final set $F$.
For every grammar symbol $X$, terminal or nonterminal, introduce fresh
\emph{nonterminal} markers $[p,X,q]$, for $p,q\in Q$.
For each grammar rule $X_1\cdots X_k\to Y_1\cdots Y_l$, introduce all rules
\[
 [p_0,X_1,p_1]\cdots[p_{k-1},X_k,p_k]
 \;\to\;
 [r_0,Y_1,r_1]\cdots[r_{l-1},Y_l,r_l],
\]
where $p_0=r_0$ and $p_k=r_l$; internal states range independently over $Q$.
All have $l\geq k$.  Include $[p,a,\delta(p,a)]\to a$ for every terminal
$a$, and fresh-start rules $T\to[q_0,S,f]$ for $f\in F$.
Include $T\to\eps$ exactly when both original languages contain $\eps$.
The new start never appears on a right-hand side.

In any new derivation, keep the endpoint annotations as ghost data even
after a terminal marker is finalized.  Adjacent endpoints agree initially
and after every rule.  A finalized terminal realizes the corresponding
one-letter automaton transition.  Projecting markers to their old symbols
turns the other rules into old grammar steps; finalization projects to an
identity step.  Hence a final terminal word belongs to the old grammar
and labels a path from $q_0$ to a final state.

Conversely, take an old successful derivation and an accepting automaton
run on its terminal yield.  Annotate the final word by the run.  Work
backwards through the derivation: the replaced right-hand segment already
has endpoints, and any choice of internal states annotates the old
left-hand segment with those same endpoints.  Thus each old step becomes
one of the enumerated annotated rules.  The beginning is
$[q_0,S,f]$ for the same final state $f$.  Run these annotated steps forwards
and finalize all terminal markers at the end.  The backwards annotation
is recursion over a supplied finite trace; its invariant is a bounded
endpoint check.  This establishes the converse and primitive recursive
witness translations.
\end{proof}

\begin{proposition}[An explicit noncontracting grammar beyond context-free
languages]\label{eff:abc}
The following noncontracting grammar generates exactly
$\{a^n b^n c^n:n\geq1\}$:
\[
\begin{gathered}
 S\to aSBC\mid aBC,\qquad CB\to BC,\\
 aB\to ab,\qquad bB\to bb,\qquad bC\to bc,\qquad cC\to cc.
\end{gathered}
\]
Adding a fresh isolated start $S_0\to S\mid\eps$ generates exactly
$\{a^n b^n c^n:n\in\N\}$ under the permitted empty-word convention.
Thus the non-context-free example established by the context-free pumping
argument occurs at the noncontracting grammar level.
\end{proposition}
\begin{proof}
Apply the recursive start rule $n-1$ times and the base rule once to obtain
$a^n(BC)^n$.  Repeatedly replace an adjacent $CB$ by $BC$.  Each replacement
reduces the number of pairs consisting of a $C$ to the left of a $B$ by
one, so at most $n^2$ replacements yield $a^n B^n C^n$.  Convert the $B$'s
from left to right by $aB\to ab$ and $bB\to bb$, then convert the $C$'s by
$bC\to bc$ and $cC\to cc$.  This gives a primitive recursive derivation
of the asserted word.

For soundness, while $S$ is present every sentential form is $a^j S u$,
where $u$ is a word in $B,C$ with equally many of each, namely $j$.
Swaps may already occur in $u$, but no terminalization rule applies because
$S$ separates its left boundary from the terminal $a$'s.  After the base
rule removes $S$, the form is $a^n u$ with $n$ copies of each of $B,C$.
Thereafter every reachable form has shape $a^n b^i c^j v$ with
$v\in\{B,C\}^*$.  The only terminalization is at the boundary of the
terminal prefix; after a $c$ has appeared no remaining $B$ can be
terminalized past it.  Swapping preserves this shape.  Counts of
$B$ plus $b$, and of $C$ plus $c$, are both invariantly $n$.
If the final word has no nonterminals, it is therefore $a^n b^n c^n$.
The fresh-start extension adds exactly the empty word: its first rule is
either the empty alternative or the one entering the preceding grammar,
and $S_0$ never occurs again.  All invariants are finite-word characteristic
checks on supplied traces.
\end{proof}

\subsection{Diagonal obstructions, without excluded middle}

Write
\[
 H(e,x)\;:\!\Longleftrightarrow\;\exists t\,
 \operatorname{Halt}(e,x,t)=1,
 \qquad K(e)\;:\!\Longleftrightarrow\;H(e,e).
\]
By Theorem~\ref{eff:three-presentations}, $K$ is recognizable.  A fixed
machine can first copy its input and then simulate the machine whose code
is that input.  An interpreter maintains the simulated table and tape
configuration, performs the displayed primitive recursive step operation,
and repeats.  Given a finite simulated trace, the fixed step compiler of
Lemma~\ref{eff:compile} supplies a primitive recursive bound for the
interpreter's finite trace; conversely each completed simulated step is
checked locally.  This is all the universal simulation needed below.

\begin{theorem}[No total Boolean decider for diagonal halting]
\label{eff:halting}
\ipra\ proves that no coded Boolean machine can satisfy both the totality
premise of Definition~\ref{eff:def-total} and the assertion
\[
 \forall e\;\bigl(K(e)\ \longleftrightarrow\
       \exists t\,\operatorname{Out}(d,e,t,1)=1\bigr).
\]
Consequently the recognizable language of canonical binary codes of $K$
is not recursive, and in particular is not primitive recursive.
\end{theorem}
\begin{proof}
From a proposed decider code $d$, construct a machine $P_d$ which, on input
$x$, simulates $d$ on $x$, halts if its Boolean output is $0$, and enters
an absorbing nonhalting loop if the output is $1$.  A non-Boolean halt may
also loop.  Inserting the fixed interpreter and these two branches is a
primitive recursive transformation on finite programs; let $p$ be its
code.  A trace showing output $0$ gives a halting trace of $P_d$ with a
primitive recursive overhead bound.  A trace showing output $1$ proves
by quantifier-free induction on any later time bound that $P_d$ never
halts: before the simulated halt it is in its interpreter states, and
afterwards it is in its designated loop state.

Apply the assumed totality to input $p$ and inspect the resulting Boolean
bit.  If it is $0$, the construction gives $H(p,p)$, hence $K(p)$; the
assumed correctness gives an output-$1$ trace of $d$ on $p$, contradicting
its output-$0$ trace.  If the bit is $1$, correctness gives $K(p)$ and a
finite halting trace for $P_d$ on $p$; the preceding nonhalting invariant
contradicts that trace.  This exhausts a decidable two-element datum,
without deciding $K(p)$ by excluded middle.  The binary-word version is
obtained by the fixed primitive recursive encoding and decoding.
\end{proof}

\begin{theorem}[The complement has no recognizer]
\label{eff:complement}
There is no machine recognizing the complement of $K$.  More precisely,
\ipra\ refutes the assertion that a given recognizer accepts input $e$
exactly when $\neg K(e)$.  Thus recognizable languages are not closed under
complement.  This conclusion does not rely on the classical form of the
two-recognizer theorem.
\end{theorem}
\begin{proof}
Given the proposed recognizer $r$, construct $Q_r$ which on input $x$
simulates $r$ on $x$, halts if $r$ accepts, and continues forever if $r$
rejects.  If $r$ itself runs forever, so does $Q_r$.  Trace translations
prove $H(Q_r,x)$ equivalent to acceptance by $r$ on $x$.
Let $q$ be the code of $Q_r$.  The proposed correctness supplies
\[
 K(q)\longleftrightarrow\neg K(q).
\]
This is constructively contradictory: the forward implication first gives
$\neg K(q)$, since assuming $K(q)$ produces both it and its negation;
the reverse implication then gives $K(q)$.  Every step is intuitionistic.
For any fixed primitive recursive verifier claimed to present $\neg K$,
Theorem~\ref{eff:three-presentations} compiles it to such a recognizer, so
the same contradiction refutes that claim.  The latter assertion is a
schema over fixed verifier definitions, not quantification over an alleged
primitive recursive universal evaluator.
\end{proof}

\begin{theorem}[Erasing a primitive recursive language can destroy
decidability]\label{eff:erasing}
There is a primitive recursive language $T$ over
$\{0,1,\#,c\}$ and a homomorphism $h$ with
$h(0)=0$, $h(1)=1$, $h(\#)=h(c)=\eps$, such that $h(T)$ is the binary-coded
diagonal halting language.  In particular, primitive recursive and recursive
languages are not closed under arbitrary erasing homomorphisms.
\end{theorem}
\begin{proof}
Let $\operatorname{bin}(e)$ be the canonical nonempty binary representation
of $e$, with $0$ represented by the single symbol $0$.  Define
\[
 T=\{\operatorname{bin}(e)\,\#\,c^t:
               \operatorname{Halt}(e,e,t)=1\}.
\]
To test membership, check the single delimiter, canonical binary prefix,
and all-$c$ suffix, decode $e$, count $t$, and apply the bounded interpreter.
All operations are primitive recursive, even though the decoded numeric
$e$ can be much larger than the input length.  If $v\in T$, then
$h(v)=\operatorname{bin}(e)$ and its displayed $t$ proves $K(e)$.
Conversely, any witness $t$ for $K(e)$ gives the explicit preimage
$\operatorname{bin}(e)\#c^t\in T$.  Therefore $h(T)=\operatorname{bin}(K)$.
Theorem~\ref{eff:halting} proves that this image has no total decider.
\end{proof}

\paragraph{Why the time parameter cannot be removed.}
The diagonal obstruction of Proposition~\ref{base:noeval} has the concrete
form $g(x)=E(x,x)+1$: for any fixed primitive recursive $E$, the function
$g$ is primitive recursive and differs from its $e$th section at the
explicit argument $e$.  Consequently the bounded machine simulator above
must not be turned into an alleged primitive recursive evaluator containing
all unary primitive recursive functions as sections.  Each compiled fixed
predicate has its own clock; no uniform primitive recursive clock for all
such programs has been asserted.

\paragraph{Scope of the hierarchy and the logical boundary.}
Finite alphabets, words, finite grammars, automata, and machine descriptions
have explicit natural-number codes.  This yields numerical enumerations
of descriptions.  We do not infer, inside first-order \ipra, a theorem
about the cardinality of the power set of all words: an arbitrary such
power set is not an object of the theory.  The diagonal theorems give the
appropriate effective obstruction directly.  Nor do the preceding
arguments prove the unrestricted-grammar/Turing-machine equivalence,
closure of context-sensitive languages under complement, general grammar
equivalence undecidability, or the Post correspondence problem.  Those
would require further constructions.  The supplied results establish the stated constructive inclusions and
specific separations among finite-state, context-free, noncontracting,
primitive-recursive, and recognizable presentations.  They do not claim
that every adjacent pair in every customary hierarchy has been separated.

\section{Counting, unary periodicity, and balanced parentheses}
\label{sec:counting}

This section proves numerical consequences of the finite constructions.
Formal power-series notation below always abbreviates a statement about
every supplied finite truncation. No completed infinite sequence is used
as an object of the arithmetic.

\begin{theorem}[Counting accepted words by a finite matrix]
\label{count:matrix}
For a DFA on $r\ge1$ states over $q$ letters, let
\[
 M_{ij}=\#\{a<q:\delta(i,a)=j\},\qquad
 u_i=[i=q_0],\qquad f_i=[i\text{ is accepting}].
\]
The number $c_n$ of accepted words of length $n$ is
$c_n=uM^nf$. All entries, powers, and counts are primitive recursive,
uniformly in the finite DFA and $n$, and $0\le c_n\le q^n$.
\end{theorem}
\begin{proof}
Let $v_{n,j}$ count words of length $n$ ending at state $j$.
At length zero its value is $u_j$. Every word of length $n+1$ has a
unique length-$n$ prefix and a final letter. Partition those prefixes
by their reached state and then count the letters leading to $j$:
\[
 v_{n+1,j}=\sum_{i<r}v_{n,i}M_{ij}.
\]
The partition is a finite bijection, so its equality follows by induction
on the finite enumeration, with equality of running counts as invariant.
Inducting on $n$ gives $v_n=uM^n$. Summing over accepting states proves
the formula. At each state the row sum of $M$ is $q$; thus
$\sum_jv_{n,j}=q^n$ by the same recurrence, and the accepting sum is
bounded by the total. This argument counts words, not accepting paths
of a nondeterministic automaton, where multiplicities could be wrong.
\end{proof}

\begin{lemma}[Finite adjugate identity]\label{count:adjugate}
For a finite square matrix $A$ over integer polynomials,
\[
 A\operatorname{adj}(A)=\operatorname{adj}(A)A=\det(A)I.
\]
The polynomial coefficients on both sides are primitive-recursive finite
expressions, and their equality has an \ipra{} proof with signed integers
represented by pairs of natural numbers.
\end{lemma}
\begin{proof}
Use the convention that the determinant of the empty matrix is $1$.
Define the determinant by its finite permutation sum with the sign given
by the parity of the inversion count. For fixed row $i$, group the
permutations by the column used in that row. Removing that row and column
identifies the corresponding summand with the signed minor, and gives
Laplace expansion. The sign changes by $(-1)^{i+j}$: moving row $i$ and
column $j$ to the first positions uses $i+j$ adjacent exchanges, and
the remaining inversion parity is that of the smaller permutation.
Each grouping and deletion has a bounded inverse on the corresponding
finite permutation table.

For a diagonal entry of $A\operatorname{adj}(A)$, this is exactly the
expansion of $\det(A)$. For an off-diagonal entry, it is the determinant
of the matrix obtained by repeating one row in place of another. Pair
each permutation with the permutation exchanging the columns assigned
to those two row positions. The product terms agree, their signs differ,
and the involution has no fixed points, so the finite signed sum is zero.
Column expansion gives the other product. Arithmetic of integer pairs
reduces the signed identities to equalities after moving negative sums
to the other side. Polynomial multiplication is finite convolution,
so the same reasoning applies coefficient by coefficient. These are
bounded rearrangements of at most $r!$ terms; induction is only on the
number of terms and the polynomial coefficient position.
\end{proof}

\begin{theorem}[Rational counting series and an integer recurrence]
\label{count:rational}
From any $r$-state DFA one can primitive-recursively compute integer
polynomials
\[
 D(z)=\det(I-zM)=1+d_1z+\cdots+d_rz^r,\qquad
 P(z)=u\operatorname{adj}(I-zM)f,\quad\deg P\le r-1,
\]
such that, coefficientwise,
\[
 D(z)\sum_{n\ge0}c_nz^n=P(z).
\]
In particular, for every $n\ge r$,
\[
 c_n+d_1c_{n-1}+\cdots+d_rc_{n-r}=0.
\]
These assertions are \ipra{} theorems about finite coefficients.
\end{theorem}
\begin{proof}
Fix a truncation degree $N$ and work modulo $z^{N+1}$, a finite ring of
$(N+1)$ signed coefficient cells. The finite geometric sum
$S_N=I+zM+\cdots+z^NM^N$ satisfies
$(I-zM)S_N=I$ modulo $z^{N+1}$, by adjacent-term cancellation.
Multiplying by the adjugate and using Lemma~\ref{count:adjugate} gives
$D(z)S_N=\operatorname{adj}(I-zM)$ in that finite ring.
Multiply by $u$ and $f$ to obtain the displayed coefficient identity
through degree $N$. The constant determinant is $1$, since the constant
matrix is $I$; each minor has at most $r-1$ linear factors. The coefficient
of degree $n\ge r$ therefore gives the recurrence. Since $N$ was an
arbitrary numerical parameter, universal introduction proves every
coefficient equation. No convergence, eigenvalues, complex numbers,
or analytic division is needed. Conversely the equation with constant
denominator coefficient $1$ computes each next $c_n$ uniquely from
earlier coefficients, by primitive recursion with a finite history.
\end{proof}

\begin{theorem}[Unary regular languages are effectively ultimately periodic]
\label{count:unary}
For an $r$-state unary DFA there are computed numbers $i<r$ and
$1\le p\le r$ such that membership of $a^n$ equals membership of
$a^{n+p}$ for every $n\ge i$. Conversely, a unary primitive-recursive
language with supplied $i\ge0,p>0$ and this periodicity property is accepted
by an explicitly constructed DFA with $i+p$ states. The converse also
allows $i=0$.
\end{theorem}
\begin{proof}
Among the states reached at times $0,\ldots,r$, select equal states at
$i<j\le r$ and set $p=j-i$. Run composition gives equal states at
$i+t$ and $j+t$ for every $t$, by induction on $t$. Lookup of their
acceptance bits proves the first assertion. For the converse, make a
chain of states $0,\ldots,i+p-1$ whose final successor is state $i$.
The initial state is $0$, and state $j$ has the supplied language's
membership bit at $a^j$. For $n<i+p$ the reached state is $n$; afterwards
it is $i+((n-i)\bmod p)$. Repeated use of the given period, formalized
by induction on the quotient, proves the correct acceptance bit.
\end{proof}

\subsection{The one-type Dyck language}

Write an opening parenthesis as an up-step and a closing parenthesis as
a down-step. For a prefix $v$ let $U(v)$ and $D(v)$ be its counts of
these letters. Define the language $\mathcal D$ by
\[
 U(w)=D(w)\quad\text{and}\quad
 \forall j\le |w|\;D(w[0:j])\le U(w[0:j]).
\]
This is a primitive-recursive predicate. Height can be handled as a
signed pair, or by the two displayed nonnegative counts.

\begin{theorem}[Unique first return and an unambiguous grammar]
\label{count:dyck}
The language $\mathcal D$ is generated by
$S\to\eps\mid\mathtt{(}S\mathtt{)}S$.
Every nonempty Dyck word has a unique decomposition
$\mathtt{(}u\mathtt{)}v$ with $u,v\in\mathcal D$.
The matching closing position and the unique parse tree are
primitive-recursively computable.
\end{theorem}
\begin{proof}
A nonempty Dyck word starts with an opening parenthesis. Choose the first
positive prefix length at which $U=D$; it exists at the full length and
is found by bounded search. The last letter of this prefix is closing.
Before that point the height is at least one, so removing the first and
last steps gives a Dyck word $u$. The suffix $v$ is also Dyck because
the removed prefix has height zero. Conversely an opening step, any
Dyck word, a closing step, and another Dyck word never have negative
height and end at zero. Inside the first block the height is positive,
so this closing position is the unique first return.

Repeat the decomposition on the two strictly shorter words to build
the parse tree. Its nonempty internal nodes consume a disjoint pair
of parentheses, so there are at most $|w|/2$ such nodes. A work list
of subintervals and a loop of at most $|w|+1$ expansions implement the
recursion. For a word with $n$ pairs every parse has $2n+1$ variable
nodes and $2n$ terminal nodes: its variable nodes form a full binary
tree with $n$ nonempty expansions. This gives an explicit
primitive-recursive bound on all relevant tree codes. Uniqueness follows
by bounded induction on the word length, since the root split and then
both smaller parses are forced. The induction predicate checks equality
of all valid parses within the numerical bound. This proves
both grammar equivalence and unambiguity for this specific grammar.
\end{proof}

\begin{theorem}[Catalan recurrence and exact enumeration]
\label{count:catalan}
Let $C_n$ count Dyck words with $n$ opening and $n$ closing parentheses.
Then
\[
 C_0=1,\qquad C_{n+1}=\sum_{i=0}^{n}C_iC_{n-i},\qquad
 (n+1)C_n=\binom{2n}{n}.
\]
All counts and the bijections proving the identities are primitive
recursive, and the identities are provable in \ipra.
\end{theorem}
\begin{proof}
The empty word is the sole object for $n=0$. In the unique decomposition
of a nonempty word, $u$ has $i$ pairs and $v$ has $n-i$ pairs. The
decomposition and its inverse are a bijection, giving the convolution.

For the exact formula assume first $n\ge1$. There are $\binom{2n}{n}$
balanced words before imposing the prefix condition: choose the positions
of the opening steps. A bad balanced word has a first prefix of height
$-1$. Exchange opening and closing steps in this prefix. The resulting
word has $n+1$ opening and $n-1$ closing steps. Conversely any such word
has final height $2$ and therefore has a first prefix of height $1$;
exchange its steps. The result is balanced with a first prefix of height
$-1$, and the two operations are inverses. Existence of these first
crossings follows from the fact that each step changes height by one,
proved by a scan with the previous height as invariant. Thus the bad
words number $\binom{2n}{n-1}$, and
\[
 C_n=\binom{2n}{n}-\binom{2n}{n-1}.
\]
Double-count pairs consisting of an $n$-element subset of a $2n$-element
set and a distinguished element of that subset: there are
$n\binom{2n}{n}$ pairs. Removing the distinguished element instead gives
an $(n-1)$-element subset with one of its $n+1$ missing elements marked,
so $n\binom{2n}{n}=(n+1)\binom{2n}{n-1}$. Multiply the preceding
difference by $n+1$ to obtain the claimed identity. For $n=0$ it is
$C_0=1$. Subset enumeration, prefix reflection, and its inverse are
finite algorithms; their inverse identities and counting equations use
bounded induction. No real asymptotic estimate enters the proof.
\end{proof}

\begin{corollary}[Balanced parentheses are context-free and nonregular]
\label{count:dycknonreg}
The Dyck language is context-free but is not accepted by any finite DFA.
\end{corollary}
\begin{proof}
The grammar above proves context-freeness. Intersect $\mathcal D$ with
the regular language consisting of a block of opening parentheses followed
by a block of closing parentheses. The result is exactly
$\{\mathtt{(}^n\mathtt{)}^n:n\ge0\}$. A DFA for $\mathcal D$ would,
by the product construction, give a DFA for that language, contradicting
the regular pumping argument in Section~\ref{sec:regular}.
\end{proof}

\clearpage
\section{A precise route to complete Lean formalization}
\label{sec:lean}

\subsection{Three separate obligations}

A semantic Lean theorem about \texttt{Nat}, lists, or automata does not by
itself prove that the theorem is derivable in iPRA. The completion plan
therefore has three separate obligations for every operation and theorem.

\begin{enumerate}
\item \textbf{Semantic mathematics.} Define the finite structures and
prove their ordinary Lean correctness statements. Finite alphabets and
state spaces can be \texttt{Fin k}; words are lists; finite relations
are Boolean tables. Quantified assertions about an arbitrary infinite
set of words are avoided in the numerical layer.
\item \textbf{Primitive-recursive representation.} Give numerical codecs,
well-formedness bits, and a fixed primitive-recursive expression for every
operation. Prove round trips for valid codes and prove that the coded
operation agrees with the semantic one. A function that Lean accepts as
total is not thereby known to be primitive recursive.
\item \textbf{Object-theory derivability.} Construct
\texttt{Prf PRAAx \(\Gamma\) \(\varphi\)} for the first-order formula
expressing the actual theorem. Every induction node must use a
quantifier-free formula. Bounded predicates are replaced by their bits
using the inspected \texttt{BoundedFormulas} infrastructure. The final
theorem has empty context and an explicit universally closed sentence.
\end{enumerate}

These tasks are individually substantial. The supplied files specify the
data, algorithms, invariants, witness bounds, and remaining proof targets;
they do not fill an axiom-shaped placeholder and count it as a theorem.
The report's proofs supply the mathematical content to be translated.

\subsection{Finite data and total conventions}

\begin{longtable}{@{}p{0.18\textwidth}p{0.37\textwidth}p{0.38\textwidth}@{}}
\toprule
Object & Semantic data & Numerical representation and check\\
\midrule\endhead
Alphabet and word & \texttt{List (Fin k)} & Base-$(k+1)$ word code,
positive digits, explicit length and validity bit.\\
DFA & Start, transition table, final bits & Positive state count, exactly
$nk$ table entries below $n$, exactly $n$ final bits.\\
$\eps$-NFA & Initial/final bit sets and tagged edges & Letter-or-empty
edge tags with checked endpoints; $2^n$ subset codes.\\
Regular expression & Inductive expression tree & Postorder node table;
children precede parents; each nonroot has one parent; all nodes reach
the root. Sharing must be unfolded before the linear Thompson bound.\\
Grammar & Tagged variables/terminals and rule list & Checked left-side
variables, right-side symbol bounds, start and finite rule indices.\\
Parse tree & Rule-labelled ordered rooted tree & Postorder table with
unique parents, valid local productions, and computed yields.\\
Derivation & Finite rewriting sequence & Forms, positions, and rule
indices; bounded adjacent-step checker.\\
PDA & Finite transition list, stack words & Input-or-empty tags,
pop symbol, pushed word, states, and complete-run certificate.\\
Machine & Finite table and configurations & State and two base-$b$ tape
halves; fixed Boolean output convention; bounded iteration.\\
Monoid & Finite multiplication table & Unit and associative table check;
letter map and accepting bit vector.\\
Polynomial & Finite signed coefficient list & Integer pairs and finite
convolution; degree and coefficient bounds.\\
\bottomrule
\end{longtable}

Malformed numerical descriptions receive explicit defaults. Theorems about
the intended semantics retain a validity premise. A rejecting default is
used for language descriptors; the machine chapter uses its stated
immediate-halt default for malformed machine codes. Those conventions
are independent and must not be conflated by an adapter.

\subsection{Representative endpoint formulas}

Here are the actual shapes of the closed arithmetic targets. Each symbol
in them denotes a fixed primitive-recursive function or a displayed
first-order formula; no extra sorts are implicit.

\paragraph{Determinization.}
\[
 \forall E,w\;\bigl(\operatorname{ENFAValid}(E)\land
 \operatorname{WordValid}(\operatorname{alphabet}(E),w)\bigr)
 \to \operatorname{Accept}_{\rm DFA}(\operatorname{Det}(E),w)
      =\operatorname{Accept}_{\eps\rm NFA}(E,w).
\]
Use the prefix invariant that each subset bit is exactly the corresponding
reachable-state bit. The invariant quantifies only over the supplied finite
state set and the processed prefix. Normalize empty-edge blocks to fewer
than $n$ edges before constructing or comparing path witnesses.

\paragraph{Equivalence of finite automata.}
For valid $A,B$ on the same alphabet, define a bit $\operatorname{Eq}(A,B)$
by the product reachability table. Prove
\[
 \operatorname{Eq}(A,B)=1\ \longleftrightarrow\
 \forall w\bigl(\operatorname{WordValid}(k,w)\to
                  \operatorname{Accept}(A,w)=\operatorname{Accept}(B,w)\bigr).
\]
The right-hand universal formula is not an induction motive. First prove
the finite fixed-point table properties. Then fix $w$ and induct through
its run with a bounded table-membership invariant. If the bit is zero,
the stored predecessor table supplies a separator of length less than
$|A||B|$. This supplies a concrete witness to inequivalence.

\paragraph{Grammar membership.}
For valid $G,w$, the target is
\[
 \operatorname{Parse}(G,w)=1\ \longleftrightarrow\
 \exists t\;\operatorname{Tree}(G,S,w,t)=1.
\]
The forward direction exhibits the parser's tree term. The reverse direction
fixes an arbitrary supplied tree and proves the table invariant by a scan
of its finite code. The CNF version strengthens the existential to a
primitive-recursively bounded one: the selected normal tree has $3|w|-1$
nodes for nonempty input and the explicit inverse-normalization bound
controls the original-grammar tree. Turn a bound on node count and labels
into a bound on its numerical code using the chosen table codec. Empty
input is a separate branch, never an instance of the formula $3|w|-1$.

\paragraph{Pumping.}
Compute the parser's chosen tree, a longest path, repeated-variable
positions, and the five cut positions. The first endpoint proves that
these cuts satisfy the bounded conditions. A second endpoint, with a
free numerical exponent $i$, verifies the tree constructed by repeating
the selected context $i$ times. This separation prevents a hidden
unbounded universal quantifier from entering the induction invariant.
The analogous regular construction uses repeated run states.

\paragraph{Noncontracting membership.}
The endpoint replaces the arbitrary derivation existential by reachability
in the finite graph on words of length at most $n$. A shortest witness has
fewer than $V=\sum_{j\le n}m^j$ steps. For example, after numbering graph
vertices below $V$, every path with at most $V$ vertices has a list code
bounded by a primitive-recursive expression in $V$; include its length
and the step-rule witnesses in that bound. No efficient space bound is
asserted for this exhaustive algorithm.

\paragraph{Halting refutation.}
The target denies the conjunction of the displayed Boolean totality formula
and correctness as a decider for $K$. The derivation uses totality at the
explicit code of the constructed machine and a case split on its Boolean
output. It does not assert that all externally total machines have
iPRA proofs of totality. The complement refutation instead produces the
intuitionistically impossible equivalence $A\leftrightarrow\neg A$.

\subsection{Dependency order and proof obligations}

\begin{longtable}{@{}p{0.10\textwidth}p{0.40\textwidth}p{0.43\textwidth}@{}}
\toprule
Stage & Deliverable & Main proof invariant or bound\\
\midrule\endhead
1 & Words, list and table codecs & Digit identities, range tests, and
round trips on valid codes.\\
2 & Bits, least witnesses, graph paths & Bounded-formula reflection;
finite counts and path deletion.\\
3 & DFA and $\eps$-NFA runs & Equal prefix states; state-set bit equality;
empty blocks shorter than state count.\\
4 & Closure and Kleene conversions & Boundary states, finite factor lists,
expression-interval tables, elimination round index.\\
5 & Decisions and minimization & Product reachability, $n^2$ pair table,
access-word and separator certificates.\\
6 & Finite monoids & Transformation-vector equality and bounded closure
among at most $n^n$ transformations.\\
7 & Grammar trees and normalization & Frontier equality, nullable ranks,
unit paths, and local rule provenance.\\
8 & CYK and direct chart & Smaller intervals or smaller insertion ranks;
bounded witness reconstruction.\\
9 & PDA and pumping & Proper stack segments, finite endpoint chains,
longest-path repeated variables, context-splice checks.\\
10 & Noncontracting and recognizable & Length-bounded graph; supplied
trace checker; explicit lists of certificates.\\
11 & Diagonal programs & Compiler simulation and input/output traces;
Boolean branch contradictions.\\
12 & Counting & Finite bijections, matrix convolution, determinant
involutions, first-crossing reflection.\\
\bottomrule
\end{longtable}

The \texttt{lean/} directory makes these obligations concrete.
\texttt{SemanticCore.lean} provides a small candidate semantic core with
ordinary proof scripts and no placeholders, but remains uncompiled.
The subject outlines specify structures, algorithms, and detailed
remaining theorem targets. \texttt{ObjectTheoryOutline.lean.txt} lists
the fixed PR symbols and the intended closed \texttt{Prf PRAAx} endpoints.
An obligation record in an outline is a specification whose value must
still be constructed; declaring that record does not establish its fields.

\subsection{Reuse from Mathlib and from ProveIt}

Mathlib's inspected documentation supplies semantic definitions and results
for DFAs, NFAs, empty-edge automata, regular expressions, and CFG derivations
\cite{MathlibDFA,MathlibNFA,MathlibEpsilon,MathlibRegex,MathlibCFG}.
These interfaces are useful for cross-checking semantics. They often use
arbitrary types or sets and therefore need finite Boolean-table adapters
before use in this arithmetic development. At inspection, the regular
expression documentation explicitly did not supply the entire
regex--automaton equivalence; this plan therefore includes both
conversions as work to implement, rather than assuming a nonexistent
version-independent theorem name.

ProveIt's inspected arithmetic infrastructure has the stronger object-level
judgment needed here. Its bounded-formula compiler is recursion on a
\emph{fixed syntax object in the metatheory}; its output is an arithmetic
term and an arithmetic derivation. This is compatible with
Proposition~\ref{base:noeval}, because it is not a universal primitive-recursive
evaluation function in the object theory. The same distinction must be
preserved when building our new program and proof compilers.

\subsection{Completion gates and exact current status}

The inspected repository uses \texttt{leanprover/lean4:v4.32.0}. Complete
formalization should be carried out in a checkout pinned to the recorded
revision, preserving that checkout's dependency lockfile. The necessary
gates are:
\begin{enumerate}
\item Typecheck every new module with \texttt{lake env lean} or the
corresponding Lake build target.
\item Produce the explicitly named first-order formula and a value of
\texttt{Prf PRAAx []} for every claimed object-theory endpoint.
\item Inspect the endpoint formulas for validity hypotheses, empty-input
branches, nonerasing hypotheses, totality data, and certificate bounds.
\item Audit proof dependencies for \texttt{sorryAx} and added axioms.
Metatheoretic tools used to construct a derivation must not be confused
with extra axioms of the object theory.
\item If desired, transport closed results with
\texttt{IPRABridge.conciseToStandard} and inspect the translated symbols.
\end{enumerate}
No Lean executable was available in the present workspace, and none of
these Lean build gates was executed. The finite reference checks in the
next section were executed. The human proofs, unchecked Lean blueprint,
and executable checks are deliberately reported with their own statuses.

\clearpage
\section{Coverage, limitations, and verification ledger}
\label{sec:audit}

\subsection{What has been established}

The report contains \textbf{100 numbered results}: theorems, propositions,
lemmas, and corollaries, including supporting finite-arithmetic lemmas.
Closely related identities are grouped into individual results rather than
counted separately. Definitions have their own numbered entries but are
not included in this count. A machine-readable catalogue with theorem
labels, titles, and chapter paths accompanies the archive.

\begin{center}
\begin{tabular}{@{}lr@{}}
\toprule
Development & Numbered results\\\midrule
Arithmetic and finite-certificate foundation & 7\\
Words and effective language algebra & 16\\
Regular languages and finite decisions & 24\\
Finite monoids and linear grammars & 4\\
Context-free grammars and pushdown automata & 26\\
Effective languages, noncontracting grammars, and limits & 16\\
Counting, unary periodicity, and Dyck words & 7\\\midrule
Total & 100\\\bottomrule
\end{tabular}
\end{center}

These are reconstructions of established mathematical results with explicit
constructive arithmetic content, rather than claims that the classical
language-theory theorems are new. The useful additional precision is in
the witness functions, finite bounds, representation conditions, and
separation between ordinary semantic proofs and iPRA derivations.

\subsection{Map from the requested source neighborhood}

The following immediate outgoing links were selected for their mathematical
relevance. The source article is the scope guide \cite{Target}; the proofs
are supplied in this report, with primary background and formalization
references in the bibliography.

\begin{longtable}{@{}p{0.31\textwidth}p{0.62\textwidth}@{}}
\toprule
Source article or immediate topic & Results established here\\
\midrule\endhead
Formal language; alphabet; string; concatenation & Numerical word codes,
free-monoid equations, cancellation, Levi factorization, word equations,
commuting words, counting and enumeration.\\
Kleene star; string homomorphism; language operations & Effective Boolean
and semiring laws, cut-based star membership, star identities, least
closure, Arden's lemma, homomorphism laws, erasing-image obstruction.\\
Regular language; regular expression; finite-state automaton & Full DFA,
NFA, $\eps$-NFA, regex and right/left-linear grammar equivalences;
closure, pumping, finite decisions, minimization, effective Nerode data.\\
Regular language: algebraic and counting properties & Finite monoid
recognition, syntactic monoid and contexts, integer coefficient recurrence,
finite-truncation rational series, unary ultimate periodicity.\\
Context-free language and context-free grammar & Parse trees, nullable and
productive variables, CNF, CYK, direct fixed-point parsing, parse counts,
closure, PDA equivalence, pumping, finiteness, nonclosure and separators.\\
BNF and EBNF & Finite expression-to-CFG compilation with explicit parse
certificate translations; the arithmetic syntax example from the source.\\
Formal grammar and Chomsky hierarchy & Checked rewriting, noncontracting
membership, constructive strict separations, and explicit $a^nb^nc^n$
grammar. Full arbitrary-grammar/TM equivalence is not claimed.\\
Turing machine; recursive and recursively enumerable language & Bounded
simulation, fixed-PR compilation, certificate/enumerator/recognizer
equivalence, totality-sensitive decisions, closure, halting and complement
diagonals.\\
CFG balanced-parenthesis example & Unique decomposition, unambiguous
grammar, nonregularity, Catalan recurrence and exact counting identity.\\
\bottomrule
\end{longtable}

\subsection{Closure matrix for results actually proved}

Here ``PR'' means a supplied primitive-recursive membership predicate;
``REG'' and ``CFL'' mean a finite automaton and a finite CFG; ``NC'' is the
noncontracting grammar presentation defined in Section~\ref{sec:effective}.
``RE'' means a certificate presentation. A check mark means an explicit
construction and proof above; ``no'' means a proved counterexample;
a dash means this report does not establish that cell. A superscript $c$
means that the constructive statement uses an explicit list of certificates.
These conventions are part of the statement, not typographical exceptions.

\begin{center}\small
\begin{tabular}{@{}p{0.30\textwidth}ccccc@{}}
\toprule
Operation & PR & REG & CFL & NC & RE\\\midrule
Union & $\checkmark$ & $\checkmark$ & $\checkmark$ & $\checkmark$ & $\checkmark$\\
Intersection & $\checkmark$ & $\checkmark$ & no & --- & $\checkmark$\\
Complement & $\checkmark$ & $\checkmark$ & no & --- & no\\
Concatenation & $\checkmark$ & $\checkmark$ & $\checkmark$ & --- & $\checkmark$\\
Star & $\checkmark$ & $\checkmark$ & $\checkmark$ & --- & $\checkmark^c$\\
Reversal & $\checkmark$ & $\checkmark$ & $\checkmark$ & $\checkmark$ & $\checkmark$\\
Arbitrary homomorphic image & no & $\checkmark$ & $\checkmark$ & --- & $\checkmark$\\
Nonerasing image & $\checkmark$ & $\checkmark$ & $\checkmark$ & --- & $\checkmark$\\
Inverse homomorphism & $\checkmark$ & $\checkmark$ & $\checkmark$ & --- & $\checkmark$\\
Substitution & no & $\checkmark$ & $\checkmark$ & --- & $\checkmark^c$\\
Intersection with REG & $\checkmark$ & $\checkmark$ & $\checkmark$ & $\checkmark$ & $\checkmark$\\
\bottomrule
\end{tabular}
\end{center}

The ``no'' for PR substitution follows already from
Theorem~\ref{eff:erasing}: substitute the finite singleton language
$\{h(a)\}$ for each source letter. Each singleton is PR, but this
substitution has the undecidable homomorphic image exhibited there.
The NC dashes are not negative closure assertions.

\subsection{Strict containments with explicit witnesses}

For the represented finite-alphabet classes the report proves
\[
 \mathrm{REG}\ \subsetneq\ \mathrm{CFL}\ \subsetneq\
 \mathrm{NC}\ \subsetneq\ \mathrm{PR\mbox{-}decided}
 \ \subseteq\ \mathrm{total\mbox{-}decided}
 \ \subsetneq\ \mathrm{RE}.
\]
This is a summary of explicit first-order statements and metatheoretic
schemas, not a second-order assertion about class objects inside iPRA.
The witness for the first separation is $\{a^nb^n\}$.
For the second it is $\{a^nb^nc^n\}$, with a primitive-recursive
separating word for every proposed CFG. For the third it is the diagonal
language of Theorem~\ref{eff:nc-pr-separation}, whose separating word
against grammar code $e$ is its canonical binary encoding.
For the last separation it is the diagonal halting language $K$.
The inclusion into total decisions uses the compiler and PR clock for
each fixed PR predicate; the next inclusion uses accepting traces.
No strictness assertion between PR-decided and total-decided is made here.

The hierarchy is presented across finite alphabets, with the actual
alphabets specified by the witness constructions. It is not a claim that
all these strict separations hold over a unary alphabet. The empty-word
exception for noncontracting grammars is included, so the CFL inclusion
does not silently discard $\eps$.

\subsection{Statements needing more than an unqualified classical slogan}

\paragraph{Finite right congruence.}
The constructive Myhill--Nerode converse is stated with a primitive-recursive
class map, representatives, and right-stability and acceptance assertions.
An abstract finite-index relation on an arbitrary predicate does not by
itself supply those algorithms. For a supplied DFA, the stronger canonical
minimization and syntactic-context conclusions are fully effective.

\paragraph{Quotients by arbitrary languages.}
Regular-by-regular quotients are constructed. The classically valid slogan
that a regular language has regular quotients by every language cannot be
used as an unconditional algorithm on arbitrary predicates. For instance,
if $R=\Sigma^*$, then $R/K$ is all words when $K$ is inhabited and empty
otherwise. For a uniformly PR family encoding bounded halting witnesses,
an unconditional uniform compiler would decide whether the family member
is inhabited, and hence decide halting. Supplied finite state-membership
information can support an appropriate constructive version; this report
does not assert such information comes for free.

\paragraph{Two recognizers.}
In classical computability theory recognizers for a language and its
complement yield a decider. Our internal theorem states the exhaustiveness
premise $L(w)\lor\neg L(w)$ explicitly. Recognition alone supplies a
finite witness only after the relevant side is known; it does not justify
an intuitionistic proof that a dovetailing search always stops. The
complement nonrecognizability theorem uses a direct contradiction
$A\leftrightarrow\neg A$ instead of that classical shortcut.

\paragraph{Finite certificate families.}
A list of unbounded positive certificates is a single input witness.
It is not silently extracted from
$\forall i<k\exists c\,R(i,c)$ with no bound on $c$.
For PR-decided languages, and for CFGs with their proved parsing bounds,
bounded selection supplies the necessary lists. For general recognizable
languages the list is part of the explicit star or substitution certificate.

\paragraph{Full linked neighborhood.}
This report does not prove general CFG equivalence, inclusion, ambiguity,
or universality undecidability; Post correspondence; Parikh's theorem;
context-sensitive complement closure; deterministic context-free or indexed
language closure classifications; full unrestricted-grammar/TM equivalence;
or results about infinite-word automata. It does not claim an iPRA
power-set cardinality theorem. Those questions require additional
reductions, automata constructions, or a different represented object.
They are not counted among the proved results.

\subsection{Executed checks and what they mean}

\begin{longtable}{@{}p{0.33\textwidth}p{0.60\textwidth}@{}}
\toprule
Artifact & Observed result\\\midrule\endhead
Automata reference checks & 64 binary two-state DFA models, 256 unary
two-state empty-edge models, 4,096 DFA-pair comparisons, and 36,721 word
comparisons; all passed. Independent path enumeration and cut tests check
the corresponding table and conversion algorithms.\\
Grammar reference checks & 37 grammars, including empty/unit cycles and
32 reproducibly seeded small grammars; 2,331 word comparisons between
direct chart saturation and CNF plus CYK; 1,068 reconstructed parse
certificates checked; all passed.\\
Words and counting checks & 433 codec round trips, 2,347 concatenation-cut
checks, 3,969 commuting-word cases, 1,260 star comparisons, 4,081 inverse
reflection checks, 256 matrix/enumeration comparisons, and 400 recurrence
checks; all passed. Catalan counts through $n=7$ agree with the formula.\\
Mathematical review & Independent reviews of the main chapters checked
bounds, empty-word cases, collection, grammar normalization, pumping,
minimization, and diagonals. The regex encoding was tightened to exclude
unintended DAG sharing in the Thompson construction.\\
Lean & Not run. Source is an explicitly uncompiled companion and blueprint;
no completed language-theory \texttt{Prf PRAAx} endpoint is claimed.\\
Repository & Read at the recorded immutable commit. Relevant declarations
and contracts inspected; ProveIt was not built in this session.\\
\bottomrule
\end{longtable}

The executable checks test genuinely different implementations where
possible: direct arbitrary-grammar saturation against normalization and
CYK, word enumeration against matrix counting, cut enumeration against
dynamic programming, and explicit paths against automaton tables.
They are regression evidence for the implementation, not exhaustive
evidence for unbounded theorem statements. The written proofs establish
the general constructions; full kernel certification remains the stated
formalization task.

\subsection{Reproduction}

The Python programs use only the standard library. From the archive root:
\begin{lstlisting}
python3 verification/regular_reference.py
python3 code/grammar_reference.py
python3 code/word_count_reference.py
python3 tools/build_pdf.py formal_languages_ipra.tex --out-dir build
\end{lstlisting}
The PDF builder uses a local TeX format and font-map setup when necessary;
on a normal TeX Live installation its \texttt{--system} option uses the
existing configuration. The LaTeX source is a genuine editable source
for the supplied report. The supplied Lean toolchain declaration matches
the inspected ProveIt revision; it is not evidence that the blueprint
has been compiled.

\clearpage
\begin{thebibliography}{99}
\raggedright

\bibitem{Target}
\emph{Formal language}, Wikipedia, and its relevant immediate mathematical
links. Source-neighborhood guide requested by the user; inspected
22 November 2026.
\url{https://en.wikipedia.org/wiki/Formal_language}.
Relevant linked entries include \emph{Regular language},
\emph{Context-free language}, \emph{Context-free grammar},
\emph{Kleene star}, \emph{Chomsky hierarchy}, \emph{Formal grammar},
\emph{Finite-state automaton}, and \emph{Recursively enumerable language}.

\bibitem{Sipser2020}
Michael Sipser, \emph{Theory of Computation}, MIT 18.404J/6.840,
Fall 2020. Instructor's lecture notes, especially lectures 1--10:
automata, regular expressions, CFGs, PDAs, pumping, decisions, and
undecidability. Primary teaching source for conventional theorem context.
\url{https://ocw.mit.edu/courses/18-404j-theory-of-computation-fall-2020/pages/lecture-notes/}.
Lecture 7, \emph{Decision Problems for Automata and Grammars}:
\url{https://ocw.mit.edu/courses/18-404j-theory-of-computation-fall-2020/78c0346bd81b6abcb2b1dde899adfc88_MIT18_404f20_lec7.pdf}.

\bibitem{Turing1936}
Alan M. Turing, ``On computable numbers, with an application to the
Entscheidungsproblem,'' \emph{Proceedings of the London Mathematical
Society}, series 2, 42 (1936--1937), 230--265; correction in 43,
544--546. Historical primary source for finite machine coding,
universal simulation, and diagonal methods; the modern halting
statements in this report are proved directly.
\url{https://theory.stanford.edu/~trevisan/cs172-07/turing36.pdf}.

\bibitem{ProveIt}
Vladimir Reshetnikov et al., \emph{ProveIt: Primitive Recursive Arithmetic}.
Inspected commit \texttt{3eb195cab4cc3d552a9eca78880c6d428f3e44f9},
22 November 2026. The source distinguishes checked object-theory libraries
from delivered research packets. See the accompanying
\texttt{evidence/repository\_snapshot.json} for the precise inspected paths.
\url{https://github.com/VladimirReshetnikov/ProveIt/tree/3eb195cab4cc3d552a9eca78880c6d428f3e44f9/Logic/PrimitiveRecursiveArithmetic}.
Relevant modules: \texttt{Core.Pairing}, \texttt{Core.SearchParameters},
\texttt{Core.Pigeonhole}, \texttt{Foundations.BoundedFormulas},
\texttt{Foundations.BoundedFixpoint}, and \texttt{IPRABridge.Standard}.

\bibitem{ProveItComplexity}
\emph{Computational complexity in intuitionistic primitive recursive
arithmetic}, research package prepared for Vladimir Reshetnikov,
20 November 2026. Inspected README, machine conventions, and foundational
chapter at the same ProveIt revision. Reused as a design comparison;
its general proposed iPRA derivations were not treated as checked imports.
\url{https://github.com/VladimirReshetnikov/ProveIt/tree/3eb195cab4cc3d552a9eca78880c6d428f3e44f9/Logic/PrimitiveRecursiveArithmetic/docs/proofs/foundations/ipra_complexity_foundations_2026-11-20}.

\bibitem{MathlibDFA}
Lean community, \emph{Mathlib.Computability.DFA}, official generated
documentation, inspected 22 November 2026. Semantic DFA definitions,
run laws, pumping and Boolean closure.
\url{https://leanprover-community.github.io/mathlib4_docs/Mathlib/Computability/DFA.html}.

\bibitem{MathlibNFA}
Lean community, \emph{Mathlib.Computability.NFA}, official generated
documentation, inspected 22 November 2026. NFA paths, subset construction,
equivalence and reversal.
\url{https://leanprover-community.github.io/mathlib4_docs/Mathlib/Computability/NFA.html}.

\bibitem{MathlibEpsilon}
Lean community, \emph{Mathlib.Computability.EpsilonNFA}, official generated
documentation, inspected 22 November 2026. Empty-edge closure,
path semantics and conversion to NFAs.
\url{https://leanprover-community.github.io/mathlib4_docs/Mathlib/Computability/EpsilonNFA.html}.

\bibitem{MathlibRegex}
Lean community, \emph{Mathlib.Computability.RegularExpressions}, official
generated documentation, inspected 22 November 2026. Regular-expression
syntax, matching and derivatives. At inspection the documentation's
TODO explicitly stated that the full automata equivalence was not yet
provided in that file.
\url{https://leanprover-community.github.io/mathlib4_docs/Mathlib/Computability/RegularExpressions.html}.

\bibitem{MathlibCFG}
Lean community, \emph{Mathlib.Computability.ContextFreeGrammar}, official
generated documentation, inspected 22 November 2026. Rule, derivation,
generation and reversal interfaces.
\url{https://leanprover-community.github.io/mathlib4_docs/Mathlib/Computability/ContextFreeGrammar.html}.

\end{thebibliography}

\end{document}
